\chapter{Dual complexes, skeleta, and retractions}
\index{dual complex}\index{skeleton!SNC}\index{retraction!skeletal}
\label{chap:skeleta}

A model contains far more algebraic information than a finite polyhedron.
For pluripotential theory, however, much of its vertical geometry is visible
on one such polyhedron. Its vertices are divisorial valuations, its faces
are monomial valuations, and vertical divisors become affine functions. The
polyhedron sits inside the Berkovich space and admits a canonical retraction.

It is worth saying at the outset what problem this solves. A Berkovich space
is compact but enormous: already the analytic affine line has four kinds of
points (\cref{thm:berk-a1-classification}), and in higher dimension no such
classification is available. Potential theory needs finite-dimensional
approximations on which functions can be differentiated, integrated, and
compared, and the complex route through charts and partitions of unity is
closed here. The non-Archimedean substitute is a directed system of finite
polyhedra \(\Delta_{\mathcal X}\), one for each SNC model, each embedded in
\(X^{\An}\) as a compact subset \(\Sk(\mathcal X)\), each carrying a
canonical retraction, and with dense union. Every object of the theory ---
model functions, psh functions, measures, energies --- is studied first on a
single \(\Delta_{\mathcal X}\), where it becomes piecewise affine real
analysis, and then transported to \(X^{\An}\) along the system.

\paragraph{Roadmap.}
\Cref{sec:skeleton-snc} constructs the dual complex. \Cref{sec:skeleton-monomial}
constructs the monomial valuation \(v_\alpha\) attached to a point
\(\alpha\) of a cell, proves it well defined, and deduces the skeleton
embedding. \Cref{sec:skeleton-affine} equips each cell with an integral
affine structure and proves it independent of every choice; this is what
makes \emph{rational}, \emph{integral slope}, and \emph{lattice length}
unambiguous later. \Cref{sec:skeleton-retraction} constructs
\(r_{\mathcal X}\), proves it continuous, and records the inequality
\(v_{r_{\mathcal X}(x)}\le v_x\) governing all later monotonicity.
\Cref{sec:skeleton-refinement} treats refinement: skeleta nest, rational
subdivisions are realized by blowups, SNC models are cofinal, and
\(X^{\An}\) is the inverse limit of its skeleta.
\Cref{sec:skeleton-model-functions} establishes the dictionary between
vertical divisors and piecewise affine functions in both directions.
\Cref{sec:essential-skeleton} distinguishes model from canonical skeleta and
says which is used later. \Cref{sec:skeleton-worked} carries out three
complete computations.

\paragraph{Dependencies.}
The reader needs the vocabulary of \cref{chap:valuations} (centers, rational
rank, the Abhyankar inequality \Cref{thm:abhyankar-relative}), the model
language of \cref{chap:models} (vertical divisors, domination, normalized
blowups, \Cref{ex:snc-chain-models}), and from \cref{chap:berkovich} the
definition of \(X^{\An}\), the center map, and
\Cref{thm:berk-proper-compact}. Conversely this chapter is used in
\cref{chap:metrics} to build model metrics, in \cref{part:bfj} as the domain
on which psh functions are equicontinuous, in \cref{part:cld} as the target
of tropicalization, and in \cref{chap:mp-dual-complexes} for the complex
analogue.

\paragraph{Standing hypotheses.}
Unless a statement says otherwise, \(K\) is complete and discretely valued,
\(v(\pi)=1\), \(R=K^\circ\), and the residue field has characteristic zero;
\(X\) is smooth, geometrically connected, and projective over \(K\). These are the
foundational BFJ hypotheses. Two do real work. Discreteness makes
\(\mathcal X_s\) a divisor with integer multiplicities, hence makes the
normalization \(v_\alpha(\pi)=1\) available and the cells compact. Residue
characteristic zero does its decisive work in exactly one place: it supplies
resolution of singularities, hence cofinality of SNC models
(\Cref{prop:snc-cofinal}). Everything concerning a \emph{single} fixed SNC
model is independent of \emph{that} input.  It is not independent of equal
characteristic zero altogether: the presentation of
\(\widehat{\mathcal O}_{\mathcal X,\xi_J}\) as a power series ring over
\(\kappa(\xi_J)\) used in \cref{sec:skeleton-monomial} needs a Cohen
coefficient field, which exists only in equal characteristic.  The
monomial valuation \(v_\alpha\) can be built instead from the graded ring of
the \(\alpha\)-weighted filtration, in any characteristic; the coefficient
field is used only for readability.

The principal geometric inputs are the resolution theorem
\Cref{thm:appB-resolution} and the deformation homotopy of
\Cref{thm:skeleton-deformation}, whose construction is described with
its source identified.  The comparison of curve metrics in
\Cref{rem:lattice-length} has a separate cited source.

\section{SNC models and their dual complexes}
\label{sec:skeleton-snc}

Let \(\mathcal X\) be a regular projective \(R\)-model of \(X\). Write
\[
                 \mathcal X_s=\sum_{i\in I}b_iE_i,
                 \qquad b_i\in\mathbf Z_{>0},
\]
with the \(E_i\) reduced and irreducible.

Two remarks on this decomposition. First, \(\mathcal X_s\) is the divisor of
the function \(\pi\), so it is a Cartier divisor, and \(b_i=\ord_{E_i}(\pi)\)
is a positive integer because \(\mathcal X\) is flat over \(R\) with
\(\mathcal X_K=X\) dense: no component of \(\mathcal X\) lies in the special
fibre. Second, the index set \(I\) is finite, since \(\mathcal X_s\) is a
noetherian scheme.

\begin{definition}[SNC model]
The model \(\mathcal X\) is an \emph{SNC model} if the support
\(\sum_iE_i\) is a simple normal crossings divisor. For notational
simplicity, we also require every nonempty intersection
\[
                         E_J=\bigcap_{j\in J}E_j
\]
to be irreducible. Without the last requirement, cells are indexed by the
irreducible components of the \(E_J\); all constructions below are
unchanged.
\end{definition}

The definition of a simple normal crossings divisor on a regular scheme is
recalled in \cref{chap:models}: every component \(E_i\) is regular, and at
every point the components through it form part of a regular system of
parameters. We record the consequence that will be used throughout.

\begin{lemma}
\label{lem:snc-normal-form}
Let \(\mathcal X\) be an SNC model, let \(J\subset I\) be nonempty with
\(E_J\neq\varnothing\), and let \(\xi_J\) denote the generic point of
\(E_J\). Then:
\begin{enumerate}
\item\label{itm:snc-nf:regular} the local ring
  \(\mathcal O_{\mathcal X,\xi_J}\) is regular of dimension \(\#J\), and any
  choice of local equations \(z_j\in\mathcal O_{\mathcal X,\xi_J}\) for the
  components \(E_j\), \(j\in J\), is a regular system of parameters;
\item\label{itm:snc-nf:pi} there is a unit
  \(u\in\mathcal O_{\mathcal X,\xi_J}^\times\) with
  \[
                         \pi=u\prod_{j\in J}z_j^{b_j};
  \]
\item\label{itm:snc-nf:strata} \(E_J\) has codimension \(\#J\) in
  \(\mathcal X\); in particular \(\#J\le n+1\), where \(n=\dim X\), and
  \(\#J=n+1\) forces \(E_J\) to be a closed point of \(\mathcal X\).
\end{enumerate}
\end{lemma}

\begin{proof}
Let \(x\in E_J\) be a closed point and let \(J(x)\supseteq J\) be the set of
all indices \(i\) with \(x\in E_i\). By the definition of simple normal
crossings there is a regular system of parameters of
\(\mathcal O_{\mathcal X,x}\) whose first \(\#J(x)\) members are local
equations \(z_i\) of the \(E_i\), \(i\in J(x)\). Consequently
\(\bigcap_{i\in J(x)}E_i\) is regular of codimension \(\#J(x)\) at \(x\),
and the same holds for the subintersection \(E_J\), which is cut out near
\(x\) by the subfamily \((z_j)_{j\in J}\) of a regular system of parameters.
This proves the codimension statement in
\Cref{lem:snc-normal-form}\,(3); the bound \(\#J\le n+1\) follows since
\(\dim\mathcal X=n+1\), with equality only when \(E_J\) is zero
dimensional. Localizing at the generic point \(\xi_J\) kills the equations
\(z_i\) with \(i\notin J\), which become units there, and leaves
\((z_j)_{j\in J}\) as a regular system of parameters of
\(\mathcal O_{\mathcal X,\xi_J}\); this is
\Cref{lem:snc-normal-form}\,(1). Two local equations of the same \(E_j\) differ
by a unit, so the choice is immaterial.

For \Cref{lem:snc-normal-form}\,(2), the divisor of \(\pi\) on \(\mathcal X\) is
\(\mathcal X_s=\sum_ib_iE_i\). Near \(\xi_J\) the only components of
\(\mathcal X_s\) that meet a small enough neighbourhood are the \(E_j\),
\(j\in J\), so \(\pi\) and \(\prod_jz_j^{b_j}\) have the same divisor there.
Their ratio is a regular function with no zeros and no poles on a
neighbourhood of \(\xi_J\) in the regular, hence normal, scheme
\(\mathcal X\), and is therefore a unit in
\(\mathcal O_{\mathcal X,\xi_J}\).
\end{proof}

Thus, near a point of \(E_J\), one can choose regular parameters \(z_j\),
\(j\in J\), defining the components, and a unit \(u\) such that
\[
                         \pi=u\prod_{j\in J}z_j^{b_j}.
\]
The whole chapter is an elaboration of this one formula. Everything
polyhedral comes from the exponents \(b_j\); everything valuative comes from
the freedom to assign arbitrary nonnegative values to the \(z_j\).

\begin{definition}[Dual complex]
For each nonempty \(E_J\), let
\[
 \sigma_J=
 \left\{(\alpha_j)_{j\in J}\in\mathbf R_{\geq0}^{J}:
                         \sum_{j\in J}b_j\alpha_j=1\right\}.
\]
If \(J'\subset J\), the face on which \(\alpha_j=0\) for
\(j\in J\setminus J'\) is identified with \(\sigma_{J'}\). The resulting
finite simplicial complex is the \emph{dual complex}
\(\Delta_{\mathcal X}\).
\end{definition}

The vertex corresponding to \(E_i\) has coordinate \(\alpha_i=b_i^{-1}\).
The incidence relation is reversed: a nonempty intersection of \(q\)
components gives a face of dimension \(q-1\).

The reversal is the source of every later dimension count. A large stratum,
such as a whole component, is a small cell, namely a vertex; a small
stratum, such as an isolated triple point, is a large cell, namely a
triangle. Blowing up makes strata smaller and cells more numerous, so the
polyhedron grows.

\begin{lemma}
\label{lem:dual-complex-simplicial}
Let \(\mathcal X\) be an SNC model of the \(n\)-dimensional variety \(X\)
and set
\[
 \Sigma_{\mathcal X}
 =\{J\subset I:\ J\neq\varnothing,\ E_J\neq\varnothing\}.
\]
Then \(\Sigma_{\mathcal X}\) is a finite abstract simplicial complex: it is
nonempty, and if \(J\in\Sigma_{\mathcal X}\) and
\(\varnothing\neq J'\subset J\), then \(J'\in\Sigma_{\mathcal X}\).
Its geometric realization is \(\Delta_{\mathcal X}\), a compact topological
space of dimension at most \(n\), covered by the finitely many simplices
\(\sigma_J\), and the face identifications of the definition are the only
ones. Moreover \(\Delta_{\mathcal X}\) is connected.
\end{lemma}

\begin{proof}
Finiteness of \(I\) has already been noted, so \(\Sigma_{\mathcal X}\) is
finite. It is nonempty because \(\mathcal X_s\neq\varnothing\): the model is
proper and flat over the local ring \(R\), so the special fibre is a
nonempty scheme, and it has at least one irreducible component. Closure
under passage to nonempty subsets is immediate, since
\(E_{J'}\supseteq E_J\neq\varnothing\) for \(J'\subset J\).

Each \(\sigma_J\) is the intersection of the affine hyperplane
\(\{\sum_jb_j\alpha_j=1\}\) with the closed orthant
\(\mathbf R_{\geq0}^J\). Since all \(b_j>0\), this is a compact simplex with
the \(\#J\) vertices \(b_j^{-1}e_j\), and the face of \(\sigma_J\) obtained
by setting \(\alpha_j=0\) for \(j\notin J'\) is precisely \(\sigma_{J'}\).
Thus the family \((\sigma_J)\) glues along faces exactly as prescribed by
the abstract complex \(\Sigma_{\mathcal X}\), and the realization is a
finite union of compact sets, hence compact.

The dimension of \(\sigma_J\) is \(\#J-1\), and
\Cref{lem:snc-normal-form}\,(3) gives \(\#J\le n+1\); hence
\(\dim\Delta_{\mathcal X}\le n\).

Finally, connectedness. The special fibre \(\mathcal X_s\) is connected:
\(\mathcal X\) is proper and flat over the henselian local ring \(R\) with
geometrically connected generic fibre \(X\), so by Zariski's connectedness
theorem the special fibre is connected as well. Suppose
\(\Delta_{\mathcal X}=A\sqcup B\) with \(A\), \(B\) nonempty, open, and
closed; being a union of cells, each of \(A\) and \(B\) is a union of
simplices, so \(I=I_A\sqcup I_B\) with no \(E_i\) (\(i\in I_A\)) meeting any
\(E_{i'}\) (\(i'\in I_B\)), for such an intersection would be a stratum
\(E_{\{i,i'\}}\neq\varnothing\) contributing an edge from \(A\) to \(B\).
Then \(\bigcup_{i\in I_A}E_i\) and \(\bigcup_{i\in I_B}E_i\) are disjoint
nonempty closed sets covering \(\mathcal X_s\), contradicting connectedness.
\end{proof}

Each face has a canonical rational affine structure. On \(\sigma_J\), its
integral affine functions are restrictions of functions
\[
                     \alpha\longmapsto\sum_{j\in J}a_j\alpha_j,
                     \qquad a_j\in\mathbf Z.
\]
The equation \(\sum b_j\alpha_j=1\) is part of this structure. It is the
reason that Euclidean coordinates chosen by drawing the simplex have no
intrinsic meaning. \Cref{sec:skeleton-affine} makes this precise and proves
that the structure depends only on the SNC model, not on any choice of
parameters, coefficient field, or embedding of the simplex into a Euclidean
space.

\begin{example}[The two smallest dual complexes]
\label{ex:dual-complex-first}
If \(\mathcal X_s\) is irreducible and reduced, the dual complex is one
point. If \(\mathcal X_s=E_0+E_1\) and the two components meet transversely
along one connected stratum, the dual complex is the interval
\[
        \{(\alpha_0,\alpha_1)\in\mathbf R_{\geq0}^2:
                          \alpha_0+\alpha_1=1\}.
\]
\end{example}

\begin{example}[Multiplicities distort the simplex]
\label{ex:dual-complex-multiplicities}
Suppose \(\mathcal X_s=E_0+2E_1\) with \(E_0\cap E_1\neq\varnothing\)
irreducible, as happens at a chart \(\pi=uz_0z_1^2\). Then
\[
  \sigma_{\{0,1\}}
  =\{(\alpha_0,\alpha_1)\in\mathbf R_{\geq0}^2:
       \alpha_0+2\alpha_1=1\},
\]
a segment whose endpoints are \((1,0)\) and \((0,\tfrac12)\). The vertex of
\(E_1\) sits at \(\alpha_1=\tfrac12\), not at \(\alpha_1=1\): the divisorial
valuation attached to a component of multiplicity \(b\) is
\(b^{-1}\ord_E\), because it must satisfy \(v(\pi)=1\). Drawing the segment
with both endpoints at Euclidean distance one from the origin would already
falsify the integral affine structure of
\Cref{sec:skeleton-affine}. Multiplicities are not cosmetic.
\end{example}

\begin{example}[A chain and a cycle]
\label{ex:dual-complex-chain-cycle}
Two one-dimensional dual complexes recur throughout the book.

\emph{Chain.} Let \(\mathcal X_s=E_0+\cdots+E_n\) with \(E_i\cap
E_{i+1}\) a single point for \(0\le i<n\) and \(E_i\cap E_j=\varnothing\)
for \(|i-j|\ge2\). Then \(\Sigma_{\mathcal X}\) consists of the \(n+1\)
singletons and the \(n\) pairs \(\{i,i+1\}\), so \(\Delta_{\mathcal X}\) is
a path with \(n+1\) vertices and \(n\) edges, homeomorphic to a segment.
This is the dual complex of the resolved conic degeneration
\(\widetilde{\mathcal C}^{(n)}\) of \Cref{ex:snc-chain-models}.

\emph{Cycle.} Let instead \(\mathcal X_s=E_0+\cdots+E_{n-1}\) with indices
in \(\mathbf Z/n\mathbf Z\), \(n\ge3\), each \(E_i\cap E_{i+1}\) a single
point and all other pairwise intersections empty. Then
\(\Delta_{\mathcal X}\) is an \(n\)-gon, homeomorphic to a circle; it is
computed in full in \Cref{ex:tate-degeneration}.

No local datum distinguishes the two: every chart of either model is the
same node chart \(z_0z_1=\pi\). The distinction is global, and it is exactly
that between a genus-zero and a genus-one degeneration.
\end{example}

\begin{example}[A triangle in dimension two]
\label{ex:dual-complex-triangle}
Let \(n=\dim X=2\) and suppose \(\mathcal X\) has a chart in which
\(\pi=xyz\) with \(x,y,z\) part of a regular system of parameters, namely
\(\mathcal U=\Spec R[x,y,z]/(xyz-\pi)\), analysed in
\Cref{ex:toric-triangle-chart}. There are three
components \(E_x,E_y,E_z\) of multiplicity one, three curve strata, and one
closed point \(E_x\cap E_y\cap E_z\); the dual complex is the full
two-simplex
\[
 \sigma=\{(\alpha_x,\alpha_y,\alpha_z)\in\mathbf R_{\geq0}^3:
          \alpha_x+\alpha_y+\alpha_z=1\}.
\]
Barycentric coordinates on it are \emph{by definition} the
\(\alpha\)-coordinates. The barycenter \((\tfrac13,\tfrac13,\tfrac13)\), the
edge midpoints such as \((\tfrac12,\tfrac12,0)\), and the vertices are the
named points at which \Cref{ex:monomial-triangle-points} evaluates the
monomial valuations.
\end{example}

\section{Monomial valuations}
\label{sec:skeleton-monomial}

The dual complex so far is a piece of combinatorics. This section attaches
to each of its points an honest valuation of the function field of \(X\),
and thereby a point of \(X^{\An}\).

\subsection{The completed local ring at a stratum}

Fix \(J\subset I\) with \(E_J\neq\varnothing\), and let \(\xi_J\) be the
generic point of \(E_J\). By \Cref{lem:snc-normal-form}\,(1) the local ring
\(\mathcal O_{\mathcal X,\xi_J}\) is regular of dimension \(\#J\), with
regular parameters \((z_j)_{j\in J}\) given by local equations of the
components. Since we are in equal characteristic zero, a coefficient
field in the completion gives
\[
 \widehat{\mathcal O}_{\mathcal X,\xi_J}
       \simeq\kappa(\xi_J)[[z_j:j\in J]].
\]
The construction can also be made intrinsically from the associated graded
ring, so it is independent of the chosen coefficient field.

Let us make both assertions precise. The ring
\(\mathcal O_{\mathcal X,\xi_J}\) is a noetherian regular local ring
containing \(\mathbf Q\), so Cohen's structure theorem provides a
coefficient field \(k_J\subset\widehat{\mathcal O}_{\mathcal X,\xi_J}\)
mapping isomorphically onto \(\kappa(\xi_J)\), and the completion is the
formal power series ring over \(k_J\) in any regular system of parameters.
The coefficient field is not unique. What is unique --- and what the
construction really uses --- is the associated graded ring for the
\(\alpha\)-weighted filtration.

Given \(\alpha\in\mathbf R_{>0}^{J}\), set for \(t\in\mathbf R\)
\[
 \mathfrak a_{\ge t}
 =\Bigl\{\textstyle\sum_\beta c_\beta z^\beta:
   c_\beta=0\text{ whenever }\langle\alpha,\beta\rangle<t\Bigr\},
\]
a decreasing family of ideals of
\(\widehat{\mathcal O}_{\mathcal X,\xi_J}\) with
\(\mathfrak a_{\ge s}\cdot\mathfrak a_{\ge t}\subset\mathfrak a_{\ge s+t}\).
Its associated graded ring is
\[
 \operatorname{gr}_\alpha
 =\bigoplus_{t\in\mathbf R}\mathfrak a_{\ge t}/\mathfrak a_{>t}
 \simeq\kappa(\xi_J)\bigl[Z_j:j\in J\bigr],
\]
a polynomial ring in which \(Z_j\) is the class of \(z_j\) in degree
\(\alpha_j\). The isomorphism is canonical once the parameters are fixed;
in particular \(\operatorname{gr}_\alpha\) is an integral domain. That is
the only property we shall use, and it is what makes the recipe below
multiplicative.

\subsection{The valuation attached to an interior point}

First suppose that \(\alpha\) lies in the relative interior of
\(\sigma_J\), so every \(\alpha_j>0\). For
\[
                     f=\sum_{\beta\in\mathbf N^J}
                          c_\beta z^\beta
\]
in the completed local ring, define
\[
 v_\alpha(f)=
 \min_{c_\beta\neq0}\langle\alpha,\beta\rangle
 =\min_{c_\beta\neq0}\sum_{j\in J}\alpha_j\beta_j.
\]
The minimum is attained: the set
\(\{\langle\alpha,\beta\rangle:\beta\in\mathbf N^J\}\) is a subset of
\(\mathbf R_{\ge0}\) that is discrete near \(0\), because
\(\langle\alpha,\beta\rangle\ge(\min_j\alpha_j)\,|\beta|\) and only finitely
many \(\beta\) satisfy \(|\beta|\le C\) for a given \(C\). This is the only
point at which positivity of all the \(\alpha_j\) is used, and it is the
reason why boundary points are treated separately below rather than by
allowing zero weights. We set \(v_\alpha(0)=+\infty\).

\begin{lemma}
\label{lem:monomial-valuation}
The function \(v_\alpha\) is a real semivaluation. Its restriction to \(K\)
is \(v\), and it induces a valuation on the function field \(K(X)\). It is
independent of the parameters up to the canonical identification determined
by the SNC pair.
\end{lemma}

\begin{proof}
We proceed in four steps.

\emph{Step 1: the axioms.} Clearly \(v_\alpha(1)=0\) and
\(v_\alpha(0)=+\infty\). For \(f=\sum c_\beta z^\beta\) and
\(g=\sum d_\beta z^\beta\) the coefficient of \(z^\beta\) in \(f+g\) is
\(c_\beta+d_\beta\), which vanishes unless \(c_\beta\neq0\) or
\(d_\beta\neq0\); hence
\[
 v_\alpha(f+g)\ \ge\ \min\{v_\alpha(f),v_\alpha(g)\},
\]
the ultrametric inequality in additive form.

\emph{Step 2: multiplicativity.} Let \(f,g\) be nonzero with
\(s=v_\alpha(f)\) and \(t=v_\alpha(g)\). By construction
\(f\in\mathfrak a_{\ge s}\setminus\mathfrak a_{>s}\), so its class
\(\operatorname{in}_\alpha(f)\in\mathfrak a_{\ge s}/\mathfrak a_{>s}\) is a
nonzero homogeneous element of \(\operatorname{gr}_\alpha\) of degree
\(s\); it is the sum of the terms \(c_\beta Z^\beta\) with
\(\langle\alpha,\beta\rangle=s\). Similarly for \(g\). The ring
\(\operatorname{gr}_\alpha\) is a polynomial ring over a field, hence a
domain, so
\(\operatorname{in}_\alpha(f)\operatorname{in}_\alpha(g)\neq0\) in degree
\(s+t\). Since this product is the class of \(fg\) in
\(\mathfrak a_{\ge s+t}/\mathfrak a_{>s+t}\), we get
\(fg\in\mathfrak a_{\ge s+t}\setminus\mathfrak a_{>s+t}\), that is,
\[
 v_\alpha(fg)=v_\alpha(f)+v_\alpha(g).
\]
In particular \(v_\alpha\) has trivial kernel on the completed local ring,
which is a domain, and therefore extends uniquely to its fraction field by
\(v_\alpha(f/g)=v_\alpha(f)-v_\alpha(g)\); the extension is well defined by
multiplicativity. Restricting along
\(K(X)\hookrightarrow\operatorname{Frac}\widehat{\mathcal O}_{\mathcal X,\xi_J}\)
--- an injection because the completion of a domain at a regular local ring
is again a domain and \(\mathcal O_{\mathcal X,\xi_J}\) has fraction field
\(K(X)\) --- yields a valuation of \(K(X)\).

\emph{Step 3: normalization on \(K\).} An element of
\(\widehat{\mathcal O}_{\mathcal X,\xi_J}^{\times}\) is a power series with
\(c_0\neq0\), so it has \(v_\alpha=0\); conversely \(v_\alpha(f)=0\) forces
\(c_0\neq0\). From \(\pi=u\prod_jz_j^{b_j}\) with \(u\) a unit, obtained in
\Cref{lem:snc-normal-form}\,(2), we get
\[
             v_\alpha(\pi)=\sum_jb_j\alpha_j=1,
\]
the last equality being the defining equation of \(\sigma_J\). Since
\(v\) is discrete with \(v(\pi)=1\) and every element of \(K^\times\) is
\(\pi^m\) times a unit of \(R\), and units of \(R\) are units of
\(\mathcal O_{\mathcal X,\xi_J}\), the restriction of \(v_\alpha\) to \(K\)
is \(v\).

\emph{Step 4: independence of the parameters.} Suppose \((z_j')_{j\in J}\)
is another system of local equations of the \(E_j\) at \(\xi_J\), and let
\(v_\alpha'\) be the resulting valuation. Write \(z_j'=u_jz_j\) with
\(u_j\in\mathcal O_{\mathcal X,\xi_J}^\times\), which is possible because
two local equations of the same prime divisor on a normal scheme differ by
a unit. Both valuations are zero on units, so
\(v_\alpha'(z_j)=v_\alpha'(z_j')=\alpha_j=v_\alpha(z_j)\). Now expand an
arbitrary \(f\) in the \(z\)-coordinates and use Step~2 twice: for the
finite truncation \(f_{<N}=\sum_{|\beta|<N}c_\beta z^\beta\) one has
\[
  v_\alpha'(f_{<N})
  \ \ge\ \min_{c_\beta\neq0,\,|\beta|<N}\langle\alpha,\beta\rangle,
\]
with equality unless two terms of minimal weight cancel to higher order,
which cannot happen because the initial forms are linearly independent over
\(\kappa(\xi_J)\) in \(\operatorname{gr}_{\alpha}'\). Since
\(f-f_{<N}\in\mathfrak m^N\) has both values at least \(N\min_j\alpha_j\),
letting \(N\to\infty\) gives \(v_\alpha'(f)=v_\alpha(f)\). Equivalently,
and this is the formulation used later, the valuation is characterized by
its values \(v_\alpha(z_j)=\alpha_j\), its center, and monomiality for the
SNC divisor; see \Cref{lem:monomial-uniqueness}.
\end{proof}

The characterization just alluded to deserves a separate statement, because
it is what makes \(v_\alpha\) an invariant of the pair
\((\mathcal X,\alpha)\) rather than of a presentation.

\begin{lemma}
\label{lem:monomial-uniqueness}
Let \(\alpha\) lie in the relative interior of \(\sigma_J\). Then
\(v_\alpha\) is the unique valuation \(w\) of \(K(X)\) with the following
three properties:
\begin{enumerate}
\item \(w\ge0\) on \(\mathcal O_{\mathcal X,\xi_J}\) and \(w>0\) on its
  maximal ideal, that is, \(w\) has center \(\xi_J\) on \(\mathcal X\);
\item \(w(z_j)=\alpha_j\) for one, equivalently any, system of local
  equations \(z_j\) of \(E_j\) at \(\xi_J\);
\item \(w\) is monomial in the \(z_j\), meaning
  \(w\bigl(\sum_\beta c_\beta z^\beta\bigr)
    =\min_{c_\beta\neq0}\langle\alpha,\beta\rangle\).
\end{enumerate}
In particular \(v_\alpha\) does not depend on the choice of coefficient
field in \(\widehat{\mathcal O}_{\mathcal X,\xi_J}\), nor on the choice of
local equations, nor on any Euclidean picture of \(\sigma_J\).
\end{lemma}

\begin{proof}
Existence with the three properties is the content of
\Cref{lem:monomial-valuation} together with the computation of the center
in \Cref{lem:skeleton-monomial-invariants} below. Uniqueness is immediate:
properties (2) and (3) determine \(w\) on
\(\widehat{\mathcal O}_{\mathcal X,\xi_J}\), and property (1) guarantees
that \(\mathcal O_{\mathcal X,\xi_J}\) maps into the valuation ring of
\(w\), so that the value on a rational function is computed from a
presentation as a ratio of elements of the local ring.

For the independence assertions, note that (1)--(3) mention no coefficient
field. Changing it replaces a chosen isomorphism
\(\widehat{\mathcal O}_{\mathcal X,\xi_J}\simeq\kappa(\xi_J)[[z_j]]\) by
another one commuting with reduction modulo the maximal ideal; the
weight filtration \(\mathfrak a_{\ge t}\), and hence \(v_\alpha\), is
defined by that filtration alone and is preserved. Independence of the
\(z_j\) is Step~4 above.
\end{proof}

\subsection{Boundary points and faces}

For a boundary point \(\alpha\in\sigma_J\), set
\[
                         J_+=\{j\in J:\alpha_j>0\}.
\]
The set \(J_+\) is nonempty, since \(\sum_jb_j\alpha_j=1\), and the point
belongs to the relative interior of the face
\(\sigma_{J_+}\). We define \(v_\alpha\) by the preceding construction in
the completed local ring at the generic point \(\xi_{J_+}\), using only the
parameters \(z_j\), \(j\in J_+\). This avoids taking a minimum in an
infinite power series with zero-weight variables, which would in general not
be attained. The following lemma shows that the recipe is consistent: the
valuation attached to \(\alpha\) does not depend on which cell of
\(\Delta_{\mathcal X}\) we regard \(\alpha\) as lying in.

\begin{lemma}
\label{lem:monomial-face-compatibility}
Let \(J'\subset J\) be nonempty sets with \(E_J\neq\varnothing\), and let
\(\alpha\in\sigma_{J'}\subset\sigma_J\), so that \(\alpha_j=0\) for
\(j\in J\setminus J'\). Then the valuation obtained by applying the
construction on the cell \(\sigma_{J'}\) coincides with the valuation
obtained by applying it on any cell \(\sigma_{J''}\) with
\(J_+\subset J''\subset J\). In particular the assignment
\(\alpha\mapsto v_\alpha\) is well defined on all of
\(\Delta_{\mathcal X}\).
\end{lemma}

\begin{proof}
Both constructions produce the valuation attached to \(J_+\) by definition,
so it suffices to prove that the construction carried out at \(\xi_{J_+}\)
is what one gets by "letting the zero weights go to zero" at
\(\xi_{J''}\); equivalently, that the valuation \(w:=v_\alpha\) defined at
\(\xi_{J_+}\) satisfies
\[
 w\Bigl(\sum_{\beta\in\mathbf N^{J''}}c_\beta z^\beta\Bigr)
 =\min_{c_\beta\neq0}\sum_{j\in J''}\alpha_j\beta_j
 =\min_{c_\beta\neq0}\sum_{j\in J_+}\alpha_j\beta_j
\]
whenever the left-hand sum is a \emph{finite} sum, that is, for
\(f\in\mathcal O_{\mathcal X,\xi_{J''}}\) written as a polynomial in the
\(z_j\), \(j\in J''\setminus J_+\), with coefficients in the local ring at
\(\xi_{J_+}\). This last statement is exactly multiplicativity plus the
ultrametric inequality for \(w\), together with \(w(z_j)=0\) for
\(j\in J''\setminus J_+\); and \(w(z_j)=0\) for such \(j\) holds because
\(z_j\) is a local equation of a component \(E_j\) not passing through
\(\xi_{J_+}\), hence a unit at \(\xi_{J_+}\).

Finally, if \(f\) is a general element of
\(\mathcal O_{\mathcal X,\xi_{J''}}\), truncate its expansion in the
\(z_j\), \(j\in J''\setminus J_+\), at order \(N\); the truncation error
lies in the \(N\)-th power of the ideal generated by those \(z_j\), on which
\(w\) is nonnegative but which is contained in
\(\mathfrak m_{\xi_{J''}}^N\). Since both sides are unchanged by truncation
once \(N\) exceeds the total degree occurring in the finitely many terms of
minimal weight, the two agree.
\end{proof}

\begin{lemma}
\label{lem:skeleton-monomial-invariants}
Let \(\alpha\in\Delta_{\mathcal X}\) and put
\(J_+=\{j:\alpha_j>0\}\), \(r=\dim_{\mathbf Q}
\bigl(\sum_{j\in J_+}\mathbf Q\alpha_j\bigr)\). Then:
\begin{enumerate}
\item the center of \(v_\alpha\) on \(\mathcal X\) is \(\xi_{J_+}\), the
  generic point of the stratum \(E_{J_+}\);
\item the value group of \(v_\alpha\) on \(K(X)^\times\) is
  \(\sum_{j\in J_+}\mathbf Z\alpha_j\), of rational rank \(r\); it contains
  \(\mathbf Z=v(K^\times)\);
\item the residue field of \(v_\alpha\) has transcendence degree
  \(n+1-r\) over the residue field \(\widetilde K\) of \(K\), and
  \(v_\alpha\) is an Abhyankar valuation, that is, equality holds in
  \Cref{thm:abhyankar-relative};
\item at a vertex, that is, when \(J_+=\{i\}\) and hence
  \(\alpha_i=b_i^{-1}\), one has
  \[
                        v_\alpha=b_i^{-1}\ord_{E_i}.
  \]
\end{enumerate}
\end{lemma}

\begin{proof}
(1) By construction \(v_\alpha\ge0\) on
\(\mathcal O_{\mathcal X,\xi_{J_+}}\) and \(v_\alpha>0\) on its maximal
ideal, since a nonunit has \(c_0=0\) and every monomial \(z^\beta\) with
\(\beta\neq0\) has weight \(\langle\alpha,\beta\rangle>0\). The center of a
valuation on a separated scheme is unique when it exists
(\Cref{prop:centers-proper-separated}), so it is \(\xi_{J_+}\).

(2) By multiplicativity the value group is generated by the values on the
local ring, and by the definition of \(v_\alpha\) every such value is of the
form \(\langle\alpha,\beta\rangle\) with \(\beta\in\mathbf N^{J_+}\).
Conversely each \(\alpha_j=v_\alpha(z_j)\) occurs. Hence the value group is
\(\sum_{j\in J_+}\mathbf Z\alpha_j\), of rational rank \(r\) by definition
of \(r\). It contains \(1=v_\alpha(\pi)\).

(3) Since \(\alpha\in\sigma_{J_+}\) has all coordinates positive and
satisfies \(\sum_{j}b_j\alpha_j=1\), we have \(1\in\sum_j\mathbf Q\alpha_j\);
therefore the rational rank of \(v_\alpha\) relative to \(v\), namely
\(\dim_{\mathbf Q}(\Gamma_{v_\alpha}/\mathbf Z)\otimes\mathbf Q\), equals
\(r-1\); note \(r\ge1\) because \(J_+\neq\varnothing\). We compute the
residue field. Let \(q=\#J_+\). The residue field of \(v_\alpha\) is the
degree-zero part of the graded fraction field of
\(\operatorname{gr}_\alpha\simeq\kappa(\xi_{J_+})[Z_j:j\in J_+]\), where
\(Z_j\) has degree \(\alpha_j\). A ratio of monomials \(Z^\beta/Z^\gamma\)
has degree zero precisely when \(\delta=\beta-\gamma\) lies in
\[
 L=\ker\bigl(\mathbf Z^{J_+}\to\mathbf R,\
   \delta\mapsto\langle\alpha,\delta\rangle\bigr),
\]
a subgroup of rank \(q-r\) by definition of \(r\). The degree-zero part is
therefore the fraction field of the group algebra
\(\kappa(\xi_{J_+})[L]\), a Laurent polynomial ring in \(q-r\) variables; so
it is a purely transcendental extension of \(\kappa(\xi_{J_+})\) of
transcendence degree \(q-r\). Since \(E_{J_+}\) has codimension \(q\) in
\(\mathcal X\) by
\Cref{lem:snc-normal-form}\,(3), its function field
\(\kappa(\xi_{J_+})\) has transcendence degree \(n+1-q\) over
\(\widetilde K\). Adding, the residue field of \(v_\alpha\) has
transcendence degree
\[
 (n+1-q)+(q-r)=n+1-r
\]
over \(\widetilde K\). Together with relative rational rank \(r-1\) this
gives \((r-1)+(n+1-r)=n\), which is exactly equality in
\Cref{thm:abhyankar-relative} for the extension \(K(X)/K\) of
transcendence degree \(n\). This proves (3) and, in particular, that the
transcendence degree over \(\widetilde K\) of the residue field of
\(v_\alpha\) is \(n+1-r\); in the divisorial case \(r=1\) it is \(n\), as it
must be.

(4) If \(J_+=\{i\}\) then \(\sigma_{\{i\}}\) is a single point with
\(b_i\alpha_i=1\), and the construction at the generic point \(\xi_i\) of
\(E_i\) reads \(v_\alpha(f)=\alpha_i\cdot\ord_{E_i}(f)\) since
\(\widehat{\mathcal O}_{\mathcal X,\xi_i}\) is a complete discrete valuation
ring with uniformizer \(z_i\). Hence \(v_\alpha=b_i^{-1}\ord_{E_i}\).
\end{proof}

\begin{remark}[Two normalizations of a divisorial point]
\label{rem:divisorial-normalization}
Part (4) explains a normalization that is easy to misremember. The vertex
of \(\Delta_{\mathcal X}\) attached to a component \(E_i\) of multiplicity
\(b_i\) is not \(\ord_{E_i}\) but \(b_i^{-1}\ord_{E_i}\). The reason is the
constraint \(v(\pi)=1\), which is forced on every point of \(X^{\An}\): a
point of the analytification of a \(K\)-variety restricts to the given
absolute value on \(K\). Over a trivially valued base this constraint
disappears and a different normalization, by log discrepancy, is used
instead; see \cref{chap:trivial-analytification}. Conflating the two is the
commonest source of factor-of-\(b\) errors in the subject.
\end{remark}

The last general fact we need about \(v_\alpha\) concerns its dependence on
\(\alpha\) for fixed \(f\). It is the tropical statement announced in the
chapter introduction, and it is the engine of every piecewise affine
argument later.

\begin{lemma}
\label{lem:monomial-newton}
Let \(J\subset I\) with \(E_J\neq\varnothing\) and let
\(f\in\mathcal O_{\mathcal X,\xi_J}\) be nonzero, with expansion
\(f=\sum_\beta c_\beta z^\beta\) in
\(\widehat{\mathcal O}_{\mathcal X,\xi_J}\). Let
\(S(f)=\{\beta:c_\beta\neq0\}\) and let
\[
 N(f)=\operatorname{conv}\bigl(S(f)+\mathbf R_{\geq0}^J\bigr)
 \subset\mathbf R_{\geq0}^J
\]
be its Newton polyhedron. Then \(N(f)\) has finitely many vertices
\(\beta^{(1)},\ldots,\beta^{(N)}\in\mathbf N^J\), all lying in \(S(f)\), and
for every \(\alpha\) in the relative interior of \(\sigma_J\),
\[
 v_\alpha(f)
 =\min_{1\le\ell\le N}\langle\alpha,\beta^{(\ell)}\rangle .
\]
Consequently \(\alpha\mapsto v_\alpha(f)\) is a concave, piecewise affine
function on \(\sigma_J\) with integral slopes, and it extends continuously
to all of \(\sigma_J\) by the same formula.
\end{lemma}

\begin{proof}
By Dickson's lemma the monomial ideal generated by \(\{z^\beta:\beta\in
S(f)\}\) is generated by finitely many of them, so \(S(f)+\mathbf N^J\) has
finitely many minimal elements and \(N(f)\) is a rational polyhedron with
finitely many vertices, each a minimal element of \(S(f)+\mathbf N^J\) and
hence in \(S(f)\).

For \(\alpha\) with all coordinates positive, the linear functional
\(\beta\mapsto\langle\alpha,\beta\rangle\) is nonnegative on
\(\mathbf R_{\geq0}^J\) and therefore attains its minimum on the polyhedron
\(N(f)\); since it is nonnegative on the recession cone, the minimum is
attained at a vertex. Because \(S(f)\subset N(f)\) and every vertex of
\(N(f)\) is in \(S(f)\), the minimum of \(\langle\alpha,\cdot\rangle\) over
\(S(f)\) equals its minimum over \(N(f)\), that is, over the vertices.

Thus \(v_\alpha(f)\) is a minimum of finitely many linear functions of
\(\alpha\) with coefficients in \(\mathbf N\), which is concave and
piecewise affine with integral slopes, and manifestly continuous on all of
\(\sigma_J\). That the extension by this formula agrees with the
definition of \(v_\alpha\) at boundary points is
\Cref{lem:monomial-face-compatibility}: setting \(\alpha_j=0\) for
\(j\notin J_+\) turns the minimum over the vertices of \(N(f)\) into the
minimum over the vertices of the image of \(N(f)\) under the coordinate
projection to \(\mathbf R^{J_+}\), which is the Newton polyhedron of \(f\)
computed at \(\xi_{J_+}\).
\end{proof}

\begin{remark}
\label{rem:newton-tropical}
\Cref{lem:monomial-newton} is the statement that valuations of regular
functions restrict to \emph{minima} of affine functions on a skeleton. Two
consequences are used constantly. First, \(-v_\alpha(f)\), which is
\(\log|f|\) in multiplicative notation, is \emph{convex} piecewise affine;
this is why vertical ideals produce convex functions in
\Cref{sec:skeleton-model-functions}, while single divisors produce affine
ones. Second, the slopes that occur are lattice points of the Newton
polyhedron, so the integral affine structure of
\Cref{sec:skeleton-affine} is exactly the structure in which these
functions are integral.
\end{remark}

\begin{theorem}\label{thm:skeleton-embedding}
There is a unique continuous injection
\[
       \iota_{\mathcal X}:\Delta_{\mathcal X}\longrightarrow X^{\mathrm{an}}
\]
whose point over \(\alpha\in\sigma_J\) is the seminorm
\(\exp(-v_\alpha)\). Its image is compact and hence closed. It is denoted
\(\operatorname{Sk}(\mathcal X)\) and called the \emph{skeleton} of
\(\mathcal X\).
\end{theorem}

\begin{proof}
\emph{Step 1: the map is well defined.} The preceding local construction
gives an analytic point on every face.
On a common subface the definitions agree by the boundary construction just
given, so the points glue. Concretely, by \Cref{lem:monomial-valuation} the
recipe produces, for each \(\alpha\in\Delta_{\mathcal X}\), a valuation
\(v_\alpha\) of \(K(X)\) restricting to \(v\) on \(K\); the associated
multiplicative seminorm \(\exp(-v_\alpha)\) on \(K(X)\) has trivial kernel
and is bounded on the local ring at its center, so it defines a point of
\(X^{\An}\) with center \(\xi_{J_+}\) by the description of analytic points
as semivaluations in \cref{sec:berk-analytification}. Compatibility across
faces is \Cref{lem:monomial-face-compatibility}.

\emph{Step 2: continuity.} The topology of \(X^{\An}\) is the coarsest one
making \(x\mapsto|f(x)|\) continuous for \(f\) a regular function on an
affine open, so it suffices to show that
\[
                   \alpha\longmapsto |f(\iota_{\mathcal X}(\alpha))|
                   =\rme^{-v_\alpha(f)}
\]
is continuous on \(\Delta_{\mathcal X}\) for every such \(f\). Since
\(\Delta_{\mathcal X}\) is a finite union of closed cells, it is enough to
prove continuity on each closed cell \(\sigma_J\); and there
\Cref{lem:monomial-newton} gives the explicit formula
\[
 v_\alpha(f)=\min_{1\le\ell\le N}\langle\alpha,\beta^{(\ell)}\rangle,
\]
a minimum of finitely many linear functions, hence continuous on all of
\(\sigma_J\) including its boundary. This settles the delicate point:
continuity up to a boundary face is not an extra argument but the statement
that the Newton-polyhedron formula remains valid when some weights vanish,
which is the last paragraph of the proof of \Cref{lem:monomial-newton}.
Applying this on each cell and using that the cells are closed and finitely
many, \(\iota_{\mathcal X}\) is continuous. The construction agrees with
that of \cite[Sections~3.1--3.2]{BFJ16}.

\emph{Step 3: injectivity.} Let \(\alpha\in\sigma_J\) and
\(\alpha'\in\sigma_{J'}\) with \(v_\alpha=v_{\alpha'}\). By
\Cref{lem:skeleton-monomial-invariants}\,(1) the two valuations have the
same center on \(\mathcal X\), so \(J_+=J'_+\); write \(J_0\) for this
common set. For \(j\in J_0\) a local equation \(z_j\) of \(E_j\) at
\(\xi_{J_0}\) satisfies \(\alpha_j=v_\alpha(z_j)=v_{\alpha'}(z_j)=\alpha'_j\),
and the remaining coordinates of both points vanish. Hence
\(\alpha=\alpha'\) as points of \(\Delta_{\mathcal X}\).

\emph{Step 4: the image.} \(\Delta_{\mathcal X}\) is compact by
\Cref{lem:dual-complex-simplicial} and \(X^{\An}\) is Hausdorff, so
\(\iota_{\mathcal X}\) is a homeomorphism onto its image and that image is
compact, hence closed. Uniqueness of a continuous map with the prescribed
values is trivial, as the values are prescribed at every point.
\end{proof}

We shall use \(\Delta_{\mathcal X}\) and \(\Sk(\mathcal X)\)
interchangeably, writing \(\alpha\) for a point of the abstract complex and
\(v_\alpha\) or \(\iota_{\mathcal X}(\alpha)\) for the corresponding
analytic point, and suppressing \(\iota_{\mathcal X}\) whenever no confusion
can arise.

The proof also supplies a useful principle: on a skeleton, valuations of
regular functions are minima of affine functions. This is the tropical
source of piecewise affine geometry.

\begin{example}[Monomial valuations at the named points of a triangle]
\label{ex:monomial-triangle-points}
Continue with the chart \(\mathcal U=\Spec R[x,y,z]/(xyz-\pi)\) of
\Cref{ex:dual-complex-triangle}, whose dual complex is the standard
triangle with barycentric coordinates
\(\alpha=(\alpha_x,\alpha_y,\alpha_z)\). All multiplicities are \(1\) and
\(v_\alpha(x)=\alpha_x\), \(v_\alpha(y)=\alpha_y\),
\(v_\alpha(z)=\alpha_z\), with \(\alpha_x+\alpha_y+\alpha_z=1\). We evaluate
at four named points, using \(f=x+y+z\) and \(g=x^2+yz\) as test functions.

\emph{The vertex \(P_x=(1,0,0)\).} By
\Cref{lem:skeleton-monomial-invariants}\,(4), \(v_{P_x}=\ord_{E_x}\), with
center the generic point of \(E_x\), value group \(\mathbf Z\), and residue
field \(\kappa(E_x)\) of transcendence degree \(2\) over \(\widetilde K\).
Here \(v_{P_x}(f)=\min\{1,0,0\}=0\) and \(v_{P_x}(g)=\min\{2,0\}=0\).

\emph{The edge midpoint \(Q=(\tfrac12,\tfrac12,0)\).} Here
\(J_+=\{x,y\}\), the center is the generic point of the curve
\(E_x\cap E_y\), the value group is \(\tfrac12\mathbf Z\), and \(r=1\), so
the residue field has transcendence degree \(2\) by
\Cref{lem:skeleton-monomial-invariants}\,(3). One computes
\(v_Q(f)=0\) and \(v_Q(g)=\min\{1,\tfrac12\}=\tfrac12\). The point is
divisorial: it is \(\tfrac12\ord_F\) for \(F\) the exceptional divisor of
the blowup of \(\mathcal U\) along \(\{x=y=0\}\), with
\(\ord_F(x)=\ord_F(y)=1\) and \(\ord_F(\pi)=2\).

\emph{The barycenter \(B=(\tfrac13,\tfrac13,\tfrac13)\).} The center is the
closed point \(E_x\cap E_y\cap E_z\), the value group is
\(\tfrac13\mathbf Z\), and \(r=1\), so the point is again divisorial: it is
\(\tfrac13\ord_F\) for \(F\) the exceptional divisor of the blowup of that
point, with \(\ord_F(x)=\ord_F(y)=\ord_F(z)=1\) and \(\ord_F(\pi)=3\). One
computes \(v_B(f)=\tfrac13\) and
\(v_B(g)=\min\{\tfrac23,\tfrac23\}=\tfrac23\); the minimum is attained
twice, which says that \(B\) lies on a break locus of the piecewise affine
function \(\alpha\mapsto v_\alpha(g)\).

\emph{An irrational interior point.} Let
\[
 \alpha=\Bigl(\tfrac{\sqrt2-1}2,\ \tfrac{\sqrt3-1}2,\
   \tfrac{4-\sqrt2-\sqrt3}2\Bigr)
 \approx(0.207,\ 0.366,\ 0.427),
\]
whose coordinates are positive and sum to \(1\). Since \(1,\sqrt2,\sqrt3\)
are \(\mathbf Q\)-linearly independent, \(r=3\), the value group has
rational rank \(3\), and by
\Cref{lem:skeleton-monomial-invariants}\,(3) the residue field is algebraic
over \(\widetilde K\). The valuation is Abhyankar of relative rational rank
\(2=\dim X\); it is quasimonomial but not divisorial, and no model of \(X\)
has a component realizing it. Such points are indispensable: a whole face,
not merely its rational points, is what makes convexity arguments and
uniform estimates work.
\end{example}

\section{The integral affine structure}
\label{sec:skeleton-affine}
\index{integral affine structure}\index{skeleton!integral affine structure}

A dual complex is not merely a topological space. Each of its cells carries
a distinguished lattice of affine functions, and virtually every later
statement --- rationality of a PL function, integrality of a slope, the
lattice length of an edge, the normalization of a Lebesgue measure on a
face --- refers to it. This section defines it, proves that it does not
depend on the choices made in \cref{sec:skeleton-monomial}, and computes it
in coordinates in the toric case.

\begin{definition}[Integral affine functions on a cell]
\label{def:integral-affine}
Let \(\sigma_J\) be a cell of \(\Delta_{\mathcal X}\). Its lattice of
\emph{integral affine functions} is
\[
 \operatorname{Aff}_{\mathbf Z}(\sigma_J)
 =\Bigl\{\alpha\longmapsto\sum_{j\in J}a_j\alpha_j\ :\
   (a_j)_{j\in J}\in\mathbf Z^{J}\Bigr\}
 \subset C^0(\sigma_J).
\]
A function is \emph{rational affine} if it lies in
\(\operatorname{Aff}_{\mathbf Z}(\sigma_J)\otimes\mathbf Q\), and \emph{affine} if it lies
in \(\operatorname{Aff}_{\mathbf Z}(\sigma_J)\otimes\mathbf R\). A \emph{rational
polyhedral subdivision} of \(\sigma_J\) is one whose cells are cut out by
finitely many inequalities \(h\ge0\) with \(h\) rational affine, and a
function is \emph{rational piecewise affine}, or rational PL, if it is
rational affine on each cell of some rational polyhedral subdivision.
\end{definition}

Three comments on the definition. First, put
\(g_J=\gcd(b_j:j\in J)\).  The subgroup of constant integral affine
functions is \(g_J^{-1}\mathbf Z\): the constant \(c\) is represented by
an integral vector precisely when \(cb_j\in\mathbf Z\) for every \(j\).
Consequently the integral-slope lattice is
\[
 \operatorname{Aff}_{\mathbf Z}(\sigma_J)/(g_J^{-1}\mathbf Z)
 \simeq \mathbf Z^J/\mathbf Z(b_j/g_J)_{j\in J},
\]
free of rank \(\#J-1=\dim\sigma_J\). Second, the representation of an
integral affine function by a vector \((a_j)\in\mathbf Z^J\) is unique:
a linear form vanishing on the affine hyperplane
\(\sum_jb_j\alpha_j=1\) must be zero. Third, the
lattice is genuinely finer than the topology: two cells can be homeomorphic
without being integral affinely isomorphic, as
\Cref{ex:dual-complex-multiplicities} already shows.

\begin{proposition}
\label{prop:integral-affine-intrinsic}
Let \(\sigma_J\) be a cell of \(\Delta_{\mathcal X}\).
\begin{enumerate}
\item For each \(j\in J\), the coordinate function \(\alpha\mapsto\alpha_j\)
  on \(\sigma_J\) equals \(\alpha\mapsto v_\alpha(z_j)\), where \(z_j\) is
  any local equation of \(E_j\) at \(\xi_J\). It is therefore independent of
  all choices.
\item Consequently \(\operatorname{Aff}_{\mathbf Z}(\sigma_J)\) depends only on the SNC
  model \(\mathcal X\) and the index set \(J\), not on the choice of local
  equations, of a coefficient field in
  \(\widehat{\mathcal O}_{\mathcal X,\xi_J}\), or of a realization of
  \(\sigma_J\) inside a Euclidean space.
\item The lattices are compatible with faces: if
  \(\varnothing\neq J'\subset J\), then restriction along
  \(\sigma_{J'}\subset\sigma_J\) maps \(\operatorname{Aff}_{\mathbf Z}(\sigma_J)\)
  \emph{onto} \(\operatorname{Aff}_{\mathbf Z}(\sigma_{J'})\). The \(\operatorname{Aff}_{\mathbf Z}\)
  therefore glue to a well-defined notion of integral affine function on
  \(\Delta_{\mathcal X}\), namely a continuous function that is integral
  affine on every cell.
\item An integral affine function on \(\sigma_J\) is the same thing as the
  restriction to \(\sigma_J\) of the model function \(\varphi_D\) of a
  vertical Cartier divisor \(D=\sum_ia_iE_i\) supported on \(\mathcal X_s\),
  in the sense of \Cref{sec:skeleton-model-functions}.
\end{enumerate}
\end{proposition}

\begin{proof}
(1) By \Cref{lem:monomial-uniqueness}\,(2), \(v_\alpha(z_j)=\alpha_j\) for
\(\alpha\) in the relative interior of \(\sigma_J\), and both sides are
continuous on \(\sigma_J\) by \Cref{lem:monomial-newton}, hence equal
everywhere. Any two local equations of \(E_j\) at \(\xi_J\) differ by a unit
of \(\mathcal O_{\mathcal X,\xi_J}\), on which every \(v_\alpha\) with
\(\alpha\in\sigma_J\) vanishes, so the function \(\alpha\mapsto
v_\alpha(z_j)\) is independent of the choice.

(2) By (1) the generators of \(\operatorname{Aff}_{\mathbf Z}(\sigma_J)\) are the
canonically defined functions \(\alpha\mapsto v_\alpha(z_j)\), and
\(v_\alpha\) itself is independent of the coefficient field by
\Cref{lem:monomial-uniqueness}. Nothing in the definition mentions a
Euclidean embedding.

(3) Restriction sends \(\sum_{j\in J}a_j\alpha_j\) to
\(\sum_{j\in J'}a_j\alpha_j\), since \(\alpha_j=0\) for
\(j\in J\setminus J'\) on \(\sigma_{J'}\); this is surjective onto
\(\operatorname{Aff}_{\mathbf Z}(\sigma_{J'})\) because the coefficients
\((a_j)_{j\in J'}\) may be prescribed arbitrarily and the remaining ones set
to zero. Gluing is then unobstructed: a family of integral affine functions,
one on each cell, agreeing on common faces, is by definition a continuous
function on \(\Delta_{\mathcal X}\) that is integral affine on every cell.

(4) This is the computation of \Cref{lem:model-function-affine}: for
\(D=\sum_ia_iE_i\) one has
\(\varphi_D(\iota_{\mathcal X}(\alpha))=\sum_{j\in J}a_j\alpha_j\) on
\(\sigma_J\). Conversely, given integers \((a_j)_{j\in J}\), extend them by
zero to an integer vector \((a_i)_{i\in I}\); the divisor
\(D=\sum_ia_iE_i\) is vertical and Cartier, because \(\mathcal X\) is
regular so that every Weil divisor is Cartier, and its model function
restricts to the prescribed integral affine function on \(\sigma_J\).
\end{proof}

\begin{remark}
\label{rem:lattice-length}
On a one-dimensional cell $\sigma_{\{i,j\}}$, integral affine
functions modulo constants have rank one.  The lattice length is
the variation of a primitive generator.  Parametrize the edge by
$t\in[0,1]$ through
\[
                 \alpha_i=t/b_i,\qquad
                 \alpha_j=(1-t)/b_j.
\]
Then
\[
 a_i\alpha_i+a_j\alpha_j
     =\frac{a_j}{b_j}
       +\frac{a_i b_j-a_j b_i}{b_i b_j}\,t .
\]
The possible numerators form $\gcd(b_i,b_j)\mathbf Z$, so the
primitive slope and lattice length are
\[
          \frac{\gcd(b_i,b_j)}{b_i b_j}
                  =\frac1{\operatorname{lcm}(b_i,b_j)}.
\]
In particular, every edge of a regular strictly semistable model
has lattice length one.

For a curve, a second useful convention is the intersection length
\[
                     \ell_{\rm int}(e)=\frac1{b_i b_j}.
\]
Write $g_e=\gcd(b_i,b_j)$.  The lattice length is
$\ell_{\rm lat}(e)=g_e\ell_{\rm int}(e)$.
These are two different metric conventions on the same graph;
the topological embedding in $X^{\An}$ does not identify them.
For algebraically closed residue field, the intersection metric
is that of \cite[(2.1.3), Theorem~2.3.3]{BN16Weight}, whereas the
lattice metric is their stable metric
\cite[Appendix~A.1.2--A.1.4]{BN16Weight}.
In residue characteristic zero, the latter agrees locally with
the geometric skeletal metric after extension to a completed
algebraic closure, away from the principal vertices
\cite[Proposition~A.2.3]{BN16Weight}.  Here a principal component
has positive genus or meets the other components in at least
three points.  This comparison is a cited input, not a consequence
of the topological retraction.

The Laplacian depends on both length and edge weight.  At a node
$q_e$ put $m_e=[\kappa(q_e):k]$.  The intersection convention uses
weight $m_e$, and the lattice convention uses weight $g_e m_e$:
\begin{equation}
\label{eq:ch05-length-weight-conversion}
 m_e\,\frac{u_j-u_i}{\ell_{\rm int}(e)}
       =g_e m_e\,\frac{u_j-u_i}{\ell_{\rm lat}(e)}.
\end{equation}
Thus either convention gives the same model measure when the
weight is changed with the length.  Unit lattice weights cannot
be retained when $g_e>1$.
For example, blowing up a node inserts multiplicity $b_i+b_j$.
The two intersection lengths add to the old one because
\[
 \frac1{b_i(b_i+b_j)}+\frac1{(b_i+b_j)b_j}
                            =\frac1{b_i b_j}.
\]
The gcd on each new edge is still $g_e$, so the same additivity
holds for lattice lengths.  Neither convention changes under
this subdivision.

For a numerical witness, assume in addition that $k$ is
algebraically closed.  On $\mathbf P^1_R$ blow up $(T,\pi)$,
then the resulting node.  The new component $F_2$ has
multiplicity two and $v_{F_2}(T)=1/2$.
Blow up a point $Q$ of $F_2$ away from its two nodes.
Its exceptional component $G$ and the strict transform of $F_2$
both have multiplicity two, and $G^2=-1$.
For $D=G$, the model function has value $1/2$ at $x_G$ and
zero at the other vertices.  Its model measure has atom
$b_G(G^2)=-2$ at $x_G$.  Both conventions reproduce it:
\[
       \frac{0-\tfrac12}{\tfrac14}
               =2\,\frac{0-\tfrac12}{\tfrac12}
               =-2 .
\]
The first expression uses intersection length $1/4$ and weight
one; the second uses lattice length $1/2$ and weight two.
The value $-1$ would result from changing the length alone.

The rational function calculation gives the same check.
Near $Q$, the residue coordinate on $F_2$ is $T^2/\pi$.
Let $\tilde c\in R^\times$ lift its nonzero value at $Q$ and put
$f=T^2-\tilde c\pi$.  Then
$v_{F_2}(f)=1$ and $v_G(f)=3/2$.
The polynomial is Eisenstein, so its zero is a closed point of
degree two; its closure meets $G$ away from the node.
The outgoing contribution at $x_G$ is again $-2$, with either
length-weight pair above.  This degree check does not select
one length convention while leaving the edge weight unspecified.
\end{remark}

The remaining question is how the lattices on two different models compare.
Since they are compared through the inclusion of one skeleton into another,
we return to it in \Cref{prop:affine-structure-refinement}, after refinement
has been set up. The answer is the expected one: the structures agree, so
that \emph{rational PL} is an unambiguous notion on
\(\bigcup_{\mathcal X}\Sk(\mathcal X)\).

\begin{example}[A toric chart and its integral affine structure in
coordinates]
\label{ex:toric-triangle-chart}
\index{toric variety!skeleton}\index{skeleton!toric}
Let \(K\) be complete discretely valued with \(v(\pi)=1\) and consider
\[
 \mathcal U=\Spec R[x,y,z]/(xyz-\pi).
\]

\emph{Regularity and strict semistability.}
The ring \(A=R[x,y,z]/(xyz-\pi)\) has dimension \(3\), being a hypersurface
in \(\Spec R[x,y,z]\). Its special fibre is
\(\Spec\widetilde K[x,y,z]/(xyz)\), the union of the three coordinate
planes, so \(\mathcal U_s=E_x+E_y+E_z\) with all multiplicities one. At the
closed point \(P=\{x=y=z=0\}\) the maximal ideal is generated by
\(x,y,z\), because \(\pi=xyz\in(x,y,z)\); the local ring has dimension
\(3\) and cotangent space spanned by \(\rd x,\rd y,\rd z\), so
\(\mathcal U\) is regular there, and at other points the same verification
uses fewer variables. Since \(\pi=xyz\) with \(x,y,z\) a regular system of
parameters, \(\mathcal U_s\) is a strictly semistable SNC divisor.

\emph{The generic fibre is a torus.} On
\(\mathcal U_K\) the product \(xyz\) equals the unit \(\pi\), so \(x,y,z\)
are invertible and \(z=\pi x^{-1}y^{-1}\); hence
\(\mathcal U_K\simeq\Spec K[x^{\pm1},y^{\pm1}]=\mathbf G_{m,K}^2\), and
\(\mathcal U\) is a toric \(R\)-model of the two-dimensional torus.

\emph{The dual complex.} The strata are the three planes, the three
coordinate axes, and \(P\), all irreducible, so \(\Delta_{\mathcal U}\) is
the full two-simplex \(\sigma\) of \Cref{ex:dual-complex-triangle}.

\emph{Tropical coordinates.} Let \(N=\mathbf Z^2\) with dual \(M\), and use
the normalization \eqref{eq:toric-trop} of
\cref{sec:toric-convex}, so that \(\Trop(w)=(v_w(x),v_w(y))\in
N_{\mathbf R}\). Under \(\iota_{\mathcal U}\) one has
\(\Trop(v_\alpha)=(\alpha_x,\alpha_y)\) with
\(\alpha_z=1-\alpha_x-\alpha_y\), so \(\Trop\) identifies \(\sigma\) with
the standard triangle
\[
 \Delta_{\mathrm{std}}
 =\{u=(u_1,u_2)\in\mathbf R_{\ge0}^2:\ u_1+u_2\le1\}
 \subset N_{\mathbf R},
\]
and \(\Sk(\mathcal U)\) is the Gauss section
\eqref{eq:toric-gauss-section} restricted to \(\Delta_{\mathrm{std}}\).

\emph{The integral affine structure, computed.} By
\Cref{def:integral-affine}, \(\operatorname{Aff}_{\mathbf Z}(\sigma)\) consists of the
functions \(a_x\alpha_x+a_y\alpha_y+a_z\alpha_z\) with
\(a_x,a_y,a_z\in\mathbf Z\). Substituting \(\alpha_z=1-u_1-u_2\),
\[
 a_x\alpha_x+a_y\alpha_y+a_z\alpha_z
 =(a_x-a_z)u_1+(a_y-a_z)u_2+a_z .
\]
As \((a_x,a_y,a_z)\) runs over \(\mathbf Z^3\), the triple
\((a_x-a_z,\ a_y-a_z,\ a_z)\) runs over all of \(\mathbf Z^3\). Hence
\[
 \operatorname{Aff}_{\mathbf Z}(\sigma)
 =\{u\longmapsto\langle m,u\rangle+c:\ m\in M,\ c\in\mathbf Z\},
\]
the standard integral affine structure of \(N_{\mathbf R}=\mathbf R^2\)
with integral constants; the elimination of \(\alpha_z\) has removed the
ambiguity noted after \Cref{def:integral-affine}.

\emph{Lattice lengths and areas.} The three edges of \(\sigma\) have
lattice length one: on the edge \(u_1+u_2=1\), for instance, a primitive
integral slope is \(u_1\), which varies from \(0\) to \(1\). The lattice
area of \(\sigma\) is \(\tfrac12\), the volume that governs total
Monge--Amp\`ere masses in \Cref{thm:toric-ma-comparison}. Neither number can
be read off a picture: both are consequences of the multiplicities
\(b_j=1\).

\emph{A distorted variant.} Replace the chart by
\(\Spec R[x,y,z]/(xy^2z^3-\pi)\), so that \(b_x=1\), \(b_y=2\),
\(b_z=3\). The dual complex is again a triangle, but now
\(\sigma=\{\alpha_x+2\alpha_y+3\alpha_z=1,\ \alpha\ge0\}\), with vertices at
\((1,0,0)\), \((0,\tfrac12,0)\), \((0,0,\tfrac13)\). Eliminating
\(\alpha_x=1-2\alpha_y-3\alpha_z\) and setting \(u=(\alpha_y,\alpha_z)\),
\[
 a_x\alpha_x+a_y\alpha_y+a_z\alpha_z
 =(a_y-2a_x)\alpha_y+(a_z-3a_x)\alpha_z+a_x,
\]
so the integral affine functions are again \(\langle m,u\rangle+c\) with
\(m\in\mathbf Z^2\), \(c\in\mathbf Z\), but the triangle now has vertices
\((0,0)\), \((\tfrac12,0)\), \((0,\tfrac13)\) in these coordinates. Its
lattice area is \(\tfrac12\cdot\tfrac12\cdot\tfrac13=\tfrac1{12}\). By
\Cref{rem:lattice-length} the three edges have lattice lengths
\[
 \frac1{\operatorname{lcm}(1,2)}=\frac12,\qquad
 \frac1{\operatorname{lcm}(1,3)}=\frac13,\qquad
 \frac1{\operatorname{lcm}(2,3)}=\frac16,
\]
which one may check directly: the last edge is \(\{\alpha_x=0\}\), running
from \((\tfrac12,0)\) to \((0,\tfrac13)\) in \(u\)-coordinates, with
direction vector \((-\tfrac12,\tfrac13)=\tfrac16(-3,2)\) and primitive
lattice direction \((-3,2)\). The three edges of a triangle in a dual
complex need not have equal lattice length, and none of this is visible in a
drawing.
\end{example}

\begin{remark}[Barycentric coordinates are not Euclidean coordinates]
\label{rem:barycentric-warning}
On the triangle of \Cref{ex:toric-triangle-chart} the integral affine
function \(h(\alpha)=\alpha_x\) decreases from \(1\) to \(0\) both along the
median from \(P_x\) to \((0,\tfrac12,\tfrac12)\) and along the edge from
\(P_x\) to \((0,1,0)\). In a Euclidean picture these two segments have
different lengths, so \(h\) would appear to have two different slopes. There
is no notion of "the" gradient of \(h\): only the lattice
\(\operatorname{Aff}_{\mathbf Z}(\sigma)\) is intrinsic, and every later
metric statement about skeleta --- edge lengths, the Laplacian on a metric
graph, Lebesgue measures on faces --- is normalized by it and by nothing
else.
\end{remark}

\section{The canonical retraction}
\label{sec:skeleton-retraction}

The skeleton is a small closed subset of a huge compact space. What makes it
useful is that the inclusion admits a canonical continuous left inverse
defined on \emph{all} of \(X^{\An}\), so that every function on the skeleton
becomes a function on the whole analytic space.

Let \(x\in X^{\mathrm{an}}\), and let \(\xi=c_{\mathcal X}(x)\). Recall from
\cref{sec:berk-reduction} that the center exists and is unique because
\(\mathcal X\) is proper over \(R\); it is the unique point of
\(\mathcal X\) at which the valuation \(v_x\) is nonnegative on the local
ring and positive on its maximal ideal. Denote by
\(J(\xi)\) the set of components \(E_j\) containing \(\xi\), and choose local
equations \(z_j\) at \(\xi\). Put
\[
                         \alpha_j=v_x(z_j).
\]
Every unit of \(\mathcal O_{\mathcal X,\xi}\) has value zero: both the unit
and its inverse belong to the valuation ring defined by \(x\). Consequently
\[
                  \sum_{j\in J(\xi)}b_j\alpha_j=v_x(\pi)=1,
\]
so \(\alpha=(\alpha_j)\) belongs to the relevant face of the dual complex.

\begin{definition}[Retraction]
The \emph{retraction associated to \(\mathcal X\)} is
\[
 r_{\mathcal X}:X^{\mathrm{an}}\longrightarrow
 \operatorname{Sk}(\mathcal X),\qquad
 r_{\mathcal X}(x)=\iota_{\mathcal X}(\alpha).
\]
\end{definition}

Three points require checking before this is a definition at all, and we
record them as a lemma.

\begin{lemma}
\label{lem:retraction-well-defined}
With notation as above:
\begin{enumerate}
\item \(J(\xi)\neq\varnothing\) and \(E_{J(\xi)}\neq\varnothing\), so
  \(\sigma_{J(\xi)}\) is a cell of \(\Delta_{\mathcal X}\);
\item \(\alpha_j=v_x(z_j)\in(0,1/b_j]\) for every \(j\in J(\xi)\), and the
  vector \(\alpha=(\alpha_j)_{j\in J(\xi)}\) lies in the relative interior
  of \(\sigma_{J(\xi)}\);
\item \(\alpha\) does not depend on the choice of local equations \(z_j\).
\end{enumerate}
\end{lemma}

\begin{proof}
(1) The center \(\xi\) lies in \(\mathcal X_s\): indeed
\(v_x(\pi)=v(\pi)=1>0\), so \(\pi\) lies in the maximal ideal of
\(\mathcal O_{\mathcal X,\xi}\), and \(\xi\) is a point of the special
fibre. Hence \(\xi\) lies on at least one component, and by construction
\(\xi\in\bigcap_{j\in J(\xi)}E_j=E_{J(\xi)}\), which is therefore nonempty.

(2) By the definition of the center, \(v_x\ge0\) on
\(\mathcal O_{\mathcal X,\xi}\) and \(v_x>0\) on its maximal ideal
\(\mathfrak m_\xi\). Since \(\xi\in E_j\), the local equation \(z_j\) lies
in \(\mathfrak m_\xi\), so \(\alpha_j=v_x(z_j)>0\); this positivity is worth
noting, because it is what makes \(J(\xi)\) rather than a proper subset the
relevant index set. Applying
\(v_x\) to the identity \(\pi=u\prod_{j\in J(\xi)}z_j^{b_j}\) of
\Cref{lem:snc-normal-form}\,(2), valid on a neighbourhood of \(\xi\)
meeting no other component, and using
\(v_x(u)=0\) gives \(\sum_jb_j\alpha_j=1\). Together with
\(\alpha_j>0\) this places \(\alpha\) in the relative interior of
\(\sigma_{J(\xi)}\) and forces \(\alpha_j\le1/b_j\).

(3) Two local equations of \(E_j\) at \(\xi\) differ by a unit of
\(\mathcal O_{\mathcal X,\xi}\), whose \(v_x\)-value is zero.
\end{proof}

Thus \(r_{\mathcal X}(x)\) records the center of \(x\) and the values of
\(x\) on the boundary components through that center, and nothing else.

\begin{proposition}
\label{prop:skeleton-retraction}
The map \(r_{\mathcal X}\) is continuous, restricts to the identity on
\(\operatorname{Sk}(\mathcal X)\), and satisfies
\[
                       r_{\mathcal X}^2=r_{\mathcal X}.
\]
If \(J(\xi)\) is the set of components through \(c_{\mathcal X}(x)\), then
the center of \(r_{\mathcal X}(x)\) is the generic point of \(E_{J(\xi)}\).
\end{proposition}

\begin{proof}
\emph{Identity on the skeleton and idempotence.} Let
\(\alpha\in\sigma_J\) and \(x=\iota_{\mathcal X}(\alpha)\). By
\Cref{lem:skeleton-monomial-invariants}\,(1), the center of \(x\) is
\(\xi_{J_+}\), whence \(J(\xi)=J_+\); and
\(v_x(z_j)=v_\alpha(z_j)=\alpha_j\) for \(j\in J_+\) by
\Cref{prop:integral-affine-intrinsic}\,(1). Thus the vector produced by the
recipe is \(\alpha\) itself, viewed in the cell \(\sigma_{J_+}\), and
\(r_{\mathcal X}(x)=x\). Since \(r_{\mathcal X}\) takes values in
\(\Sk(\mathcal X)\), idempotence \(r_{\mathcal X}^2=r_{\mathcal X}\) is an
immediate consequence.

\emph{The center.} Let \(x\in X^{\An}\) with center \(\xi\) and
\(\alpha=(v_x(z_j))_{j\in J(\xi)}\). By
\Cref{lem:retraction-well-defined}\,(2) every coordinate of \(\alpha\) is
strictly positive, so in the notation of
\Cref{lem:skeleton-monomial-invariants} we have \(J_+=J(\xi)\); that lemma,
part (1), then says that the center of
\(r_{\mathcal X}(x)=\iota_{\mathcal X}(\alpha)\) is the generic point
\(\xi_{J(\xi)}\) of \(E_{J(\xi)}\), as asserted. Note the geometric
content: the retraction moves the center from an arbitrary point \(\xi\) of
a stratum to the generic point of the smallest stratum containing it. In
particular \(r_{\mathcal X}(x)\) depends on \(\xi\) only through
\(J(\xi)\).

\emph{Continuity.} Cover the quasi-compact special fibre by finitely many
affine opens \(\mathcal V_\ell\) with the following property: the components
meeting \(\mathcal V_\ell\) have nonempty common intersection and admit
local equations \(z_{\ell j}\) on \(\mathcal V_\ell\).  Such opens are
obtained by shrinking around each point and then taking a finite subcover.
Anticontinuity of reduction
(\Cref{prop:berk-reduction-anticontinuous}) says that
\[
 \mathcal V_\ell^{\beth}
 :=c_{\mathcal X}^{-1}((\mathcal V_\ell)_s)
\]
is closed, not open.  These finitely many closed sets cover \(X^{\An}\).
On \(\mathcal V_\ell^{\beth}\) one has
\[
 r_{\mathcal X}(x)=\iota_{\mathcal X}
   \bigl((v_x(z_{\ell j}))_j\bigr),
\]
with a zero coordinate when the center misses the corresponding component.
Every \(z_{\ell j}\) is invertible on the generic fibre, so
\(x\mapsto v_x(z_{\ell j})=-\log|z_{\ell j}(x)|\) is finite and continuous
on the tube.  Hence the displayed formula, followed by
\Cref{thm:skeleton-embedding}, makes \(r_{\mathcal X}\) continuous on each
closed member of the cover.  The formulas agree on overlaps because local
equations differ by units.  The finite closed-cover gluing lemma now proves
continuity on all of \(X^{\An}\).
\end{proof}

The retraction preserves precisely the monomial data of \(x\) relative to
the boundary \(\mathcal X_s\); it forgets residue directions transverse to
that data.  The face--stratum reversal and the endpoint of the retraction are
displayed together in \cref{fig:skeleton-retraction}.

The following inequality is short to prove and is used more often than any
other statement of this chapter. It says that retracting can only make
functions bigger in absolute value, that is, valuations smaller.

\begin{lemma}
\label{lem:retraction-minorant}
Let \(x\in X^{\An}\) with center \(\xi\) on \(\mathcal X\), and let
\(f\in\mathcal O_{\mathcal X,\xi}\). Then
\[
 v_{r_{\mathcal X}(x)}(f)\ \le\ v_x(f),
 \qquad\text{equivalently}\qquad
 |f(x)|\ \le\ |f(r_{\mathcal X}(x))| .
\]
Equality holds when \(f\) is a monomial in local equations of the components
through \(\xi\), times a unit.
\end{lemma}

\begin{proof}
Write \(A=\mathcal O_{\mathcal X,\xi}\), \(J_0=J(\xi)\),
\(\alpha_j=v_x(z_j)\), \(J_+=\{j\in J_0:\alpha_j>0\}\), and
\(w=v_{r_{\mathcal X}(x)}\). We may assume \(f\neq0\); set \(s=w(f)\), a
finite nonnegative real number. If \(s=0\) the inequality is the statement
\(v_x(f)\ge0\), which holds because \(\xi\) is the center of \(x\); so
assume \(s>0\). Let \(\mathfrak q=(z_j:j\in J_+)\subset A\), a prime ideal
with \(A_{\mathfrak q}=\mathcal O_{\mathcal X,\xi_{J_+}}\), and let
\[
 \mathfrak a=\bigl(z^\beta:\ \beta\in\mathbf N^{J_+},\
   \langle\alpha,\beta\rangle\ge s\bigr)\subset A,
\]
which is finitely generated: the set of such \(\beta\) is an upward-closed
subset of \(\mathbf N^{J_+}\) and therefore has finitely many minimal
elements, by Dickson's lemma.

\emph{Step 1: \(v_x\ge s\) on \(\mathfrak a\).} Each generator satisfies
\(v_x(z^\beta)=\sum_j\beta_j\alpha_j=\langle\alpha,\beta\rangle\ge s\) by
multiplicativity, and every element of \(A\) has \(v_x\ge0\). The
ultrametric inequality applied to a finite \(A\)-linear combination of the
generators gives the claim.

\emph{Step 2: \(f\in\mathfrak a\).} This is where the definition of \(w\)
is used. Since \(f\in A\subset A_{\mathfrak q}\), expand \(f\) in the
completion
\(\widehat{A_{\mathfrak q}}\simeq\kappa(\mathfrak q)[[z_j:j\in J_+]]\);
by definition \(w(f)=s\) means that every monomial \(z^\beta\) occurring in
this expansion satisfies \(\langle\alpha,\beta\rangle\ge s\). Every such
monomial is divisible by one of the finitely many minimal generators of
\(\mathfrak a\); grouping the terms of the power series according to a
choice of divisor among those generators exhibits \(f\) as an
\(\widehat{A_{\mathfrak q}}\)-linear combination of them, that is,
\(f\in\mathfrak a\widehat{A_{\mathfrak q}}\). Since \(A_{\mathfrak q}\) is
noetherian local and its completion is faithfully flat over it,
\(\mathfrak a\widehat{A_{\mathfrak q}}\cap A_{\mathfrak q}
 =\mathfrak aA_{\mathfrak q}\), so \(f\in\mathfrak aA_{\mathfrak q}\).

It remains to descend from \(A_{\mathfrak q}\) to \(A\). The elements
\((z_j)_{j\in J_+}\) form part of a regular system of parameters of the
regular local ring \(A\), so they are a regular sequence and
\(\operatorname{gr}_{\mathfrak q}(A)\simeq(A/\mathfrak q)[Z_j:j\in J_+]\).
As \(\mathfrak a\) is a monomial ideal in these variables containing
\(\mathfrak q^M\) for \(M>s/\min_{j\in J_+}\alpha_j\), the induced
filtration of \(A/\mathfrak a\) is finite with graded pieces free over
\(A/\mathfrak q\) on the monomials \(z^\beta\notin\mathfrak a\); hence
\(A/\mathfrak a\) is a free \(A/\mathfrak q\)-module with basis
\(\{z^\beta:\langle\alpha,\beta\rangle<s\}\). A free module over the domain
\(A/\mathfrak q\) is torsion free, so
\(A/\mathfrak a\to(A/\mathfrak a)_{\mathfrak q}\) is injective, that is,
\(\mathfrak aA_{\mathfrak q}\cap A=\mathfrak a\). Therefore
\(f\in\mathfrak a\).

Combining the two steps, \(v_x(f)\ge s=w(f)\).

For the equality case, if \(f=u\prod_{j\in J_0}z_j^{\beta_j}\) with \(u\) a
unit at \(\xi\), then both sides equal
\(\sum_{j}\beta_j\alpha_j\), because \(v_x\) and \(w\) both vanish on \(u\)
and both send \(z_j\) to \(\alpha_j\).
\end{proof}

\begin{corollary}
\label{cor:retraction-ideal-monotone}
Let \(\mathfrak a\) be a coherent ideal sheaf on \(\mathcal X\) whose
zero locus is contained in \(\mathcal X_s\). Then, with
\(\log|\mathfrak a|\) defined as in
\Cref{sec:skeleton-model-functions},
\[
 \log|\mathfrak a|(x)\ \le\ \log|\mathfrak a|(r_{\mathcal X}(x))
 \qquad(x\in X^{\An}).
\]
\end{corollary}

\begin{proof}
Let \(f_1,\ldots,f_N\) generate \(\mathfrak a\) near the center \(\xi\) of
\(x\). Then
\(\log|\mathfrak a|(x)=-\min_\ell v_x(f_\ell)\) and
\(\log|\mathfrak a|(r_{\mathcal X}(x))=-\min_\ell
v_{r_{\mathcal X}(x)}(f_\ell)\), the same generators being usable at the
center of \(r_{\mathcal X}(x)\) because that center is a generalization of
\(\xi\). Apply \Cref{lem:retraction-minorant} to each \(f_\ell\) and take
minima.
\end{proof}

\begin{example}[The retraction of the standard annulus]
\label{ex:retraction-annulus}
Take the node chart \(\mathcal U=\Spec R[x,y]/(xy-\pi)\) of
\Cref{ex:snc-chain-models} with \(n=1\), whose generic fibre is the standard
annulus \(\{|\pi|\le|x|\le1\}\) of \Cref{ex:berk-standard-annulus}, with
\(y=\pi/x\). Its dual complex is the segment parametrized by
\(t=\alpha_x\in[0,1]\), and \(r_{\mathcal U}(w)\) is the point with
\(t=v_w(x)\).

Comparing with \Cref{thm:berk-a1-classification}: a Gauss point
\(\zeta_{0,\rme^{-t}}\) with \(t\in[0,1]\) is its own retraction; a type-I
point \(a\) with \(v(a)=t\in(0,1)\) retracts to \(t\), and the whole open
disc \(D^-(a,\rme^{-t})\) attached to the interval there is contracted to
that parameter, because \(|x|\) is constant on it; every point with
\(|x(w)|=1\) retracts to the vertex \(x_{E_y}\), so the preimage under
reduction of the smooth locus of \(E_y\) is crushed to one point.

This is the general picture: the retraction crushes each residue class of
the reduction map to a single skeletal point and moves points only in the
directions transverse to the skeleton. Compare \Cref{prop:berk-a1-tree},
where the interval is the unique path between \(\zeta_{0,1}\) and
\(\zeta_{0,|\pi|}\).
\end{example}

\begin{figure}[htbp]
\centering
\begin{tikzpicture}[x=1cm,y=1cm,>=Stealth,
  vertex/.style={circle,draw,fill=white,inner sep=1.8pt},
  ambientpoint/.style={circle,fill=black,inner sep=1.5pt}]
  \draw[thick,rounded corners=13pt] (-4.55,-0.75) rectangle (4.55,3.35);
  \node[anchor=west,font=\small] at (-4.25,3.02) {$X^{\mathrm{an}}$};

  \coordinate (e0) at (-2.35,0);
  \coordinate (e1) at (2.35,0);
  \coordinate (e2) at (0,2.45);
  \draw[very thick] (e0)--(e1)--(e2)--cycle;
  \node[vertex,label={[font=\small]below:$x_{E_0}$}] at (e0) {};
  \node[vertex,label={[font=\small]below:$x_{E_1}$}] at (e1) {};
  \node[vertex,label={[font=\small]right:$x_{E_2}$}] at (e2) {};
  \node[font=\scriptsize,fill=white,inner sep=1pt] at (0,-0.02)
    {$E_0\cap E_1$};
  \node[font=\scriptsize,align=center,fill=white,inner sep=1pt] at (0,1.55)
    {$E_0\cap E_1\cap E_2$\\$\longleftrightarrow$ interior};

  \coordinate (alpha) at (0.45,0.66);
  \node[ambientpoint,label={[font=\small]below right:$v_\alpha$}] at (alpha) {};
  \node[ambientpoint,label={[font=\small]right:$x$}] (x) at (3.55,2.55) {};
  \draw[-{Stealth[length=2.3mm]},thick]
    (x) .. controls (3.1,1.05) and (1.35,0.66) ..
    node[pos=0.46,below,font=\small] {$r_{\mathcal X}$} (alpha);

  \coordinate (beta) at (-1.18,1.23);
  \node[ambientpoint] at (beta) {};
  \node[ambientpoint,label={[font=\small]left:$y$}] (y) at (-3.55,1.98) {};
  \draw[-{Stealth[length=2.3mm]},thick]
    (y) to[bend right=8] node[midway,above,font=\small] {$r_{\mathcal X}$} (beta);

  \node[font=\small,fill=white,inner sep=2pt] at (0,3.05)
    {$\operatorname{Sk}(\mathcal X)\simeq\Delta_{\mathcal X}$};
\end{tikzpicture}
\caption{The displayed case \(E_0\cap E_1\cap E_2\ne\varnothing\).
Components give vertices, pairwise strata give edges, and the triple stratum
gives the two-dimensional face.  The retraction sends an analytic point to
the monomial valuation with the same boundary values
\(\alpha_i=v_x(z_i)\); the rounded rectangle is only a schematic ambient
space.}
\label{fig:skeleton-retraction}
\end{figure}

\subsection{The deformation retraction}

\Cref{prop:skeleton-retraction} produces a continuous idempotent map onto
the skeleton. In favourable circumstances much more is true: the retraction
is the endpoint of a homotopy contracting \(X^{\An}\) onto
\(\Sk(\mathcal X)\) through maps fixing the skeleton pointwise. This is the
source of every topological statement about Berkovich spaces --- local
contractibility, the homotopy type of a degeneration, the identification of
\(H^*(X^{\An})\) with the cohomology of a finite complex.

We state the theorem in Berkovich's generality and record what the rest of
the book uses. Note the hypothesis: strict polystability, in the simplest
case strict semistability with all \(b_i=1\). An SNC model with
multiplicities is \emph{not} covered, and we never assert a homotopy in that
generality.

\begin{theorem}\label{thm:skeleton-deformation}
Let \(\mathfrak X\) be a proper strictly polystable formal
\(K^\circ\)-scheme, and let \(X\) be its generic fiber. The Berkovich
skeleton of \(\mathfrak X\) is the image of a strong deformation retraction
of \(X\).
\end{theorem}
\begin{proof}[Foundational deformation input]
This is Berkovich, \emph{Smooth $p$-adic analytic spaces are locally
contractible}, Invent. Math. \textbf{137} (1999), Theorem~8.1.  Applied to
the proper strictly polystable formal scheme \(\mathfrak X\), that theorem
constructs a homotopy \(H:X\times[0,1]\to X\) whose time-one map is the
skeletal retraction and which fixes the skeleton at every time.  This is
exactly a strong deformation retraction.  The theorem is not invoked for an
SNC model with nontrivial multiplicities; there we use only the endpoint map
of \Cref{prop:skeleton-retraction}.
\end{proof}

\begin{remark}[What the book uses, and what it does not]
\label{rem:deformation-scope}
The theorem is often invoked in a stronger form than pluripotential theory
requires, so we state the dependency.
\begin{enumerate}
\item \emph{Used throughout:} the endpoint map \(r_{\mathcal X}\), its
  continuity, its idempotence, and \Cref{lem:retraction-minorant}. These
  hold for every SNC model with arbitrary multiplicities, by
  \Cref{prop:skeleton-retraction}, and their proofs do not use
  polystability.
\item \emph{Used only for orientation:} the homotopy \(H\), hence the
  assertions that \(X^{\An}\) is homotopy equivalent to
  \(\Delta_{\mathcal X}\) and has its cohomology type.  Local
  contractibility is a separate conclusion of Berkovich's theorem, not a
  formal consequence of the existence of this deformation retraction. No
  theorem proved here depends on these facts.
\item \emph{Not asserted:} a strong deformation retraction onto the skeleton
  of an SNC model with some \(b_i>1\), or of a dlt model.
\end{enumerate}
The distinction matters because the passage to the inverse system in
\Cref{thm:skeleton-inverse-limit} is a statement about \emph{endpoints
only}; homotopies do not survive that limit in a useful form.
\end{remark}

\section{Refinement, subdivision, and cofinality}
\label{sec:skeleton-refinement}

A single skeleton is a finite polyhedron and cannot see all of \(X^{\An}\).
The theory works because skeleta \emph{nest}: passing to a higher model
enlarges the polyhedron and subdivides its cells, compatibly with the
retractions. This section proves these statements, shows that rational
subdivisions are realized by models, establishes cofinality of SNC models,
and concludes that \(X^{\An}\) is the inverse limit of its skeleta.

\subsection{Nesting and compatibility of retractions}

Suppose that \(\mathcal X'\geq\mathcal X\) are SNC models in the BFJ
ordering, so \(\mathcal X'\to\mathcal X\) is a vertical blowup and is the
identity on the generic fiber. Then
\[
              \operatorname{Sk}(\mathcal X)
              \subset \operatorname{Sk}(\mathcal X').
\]
This inclusion, and the identities below, are
\cite[Proposition~3.4]{BFJ16}; they are not being asserted for an arbitrary
map between unrelated log models.
The retractions satisfy
\[
\begin{split}
 r_{\mathcal X}\circ r_{\mathcal X'}&=r_{\mathcal X},\\
 r_{\mathcal X'}\circ r_{\mathcal X}&=r_{\mathcal X},\\
 r_{\mathcal X'}|_{\operatorname{Sk}(\mathcal X)}
       &=\operatorname{id}.
\end{split}
\]
Thus \(r_{\mathcal X}|_{\operatorname{Sk}(\mathcal X')}\) is the transition
map from the finer skeleton to the coarser one.

Two of the four assertions are elementary once the inclusion is known, and
we prove them, both because the argument is instructive and because the
reader should see which one carries the content.

\begin{proposition}
\label{prop:skeleton-nesting}
Let \(\mathcal X'\geq\mathcal X\) be SNC models and assume
\(\Sk(\mathcal X)\subset\Sk(\mathcal X')\). Then
\[
 r_{\mathcal X}\circ r_{\mathcal X'}=r_{\mathcal X},
 \qquad
 r_{\mathcal X'}|_{\Sk(\mathcal X)}=\operatorname{id},
 \qquad
 r_{\mathcal X'}\circ r_{\mathcal X}=r_{\mathcal X}.
\]
\end{proposition}

\begin{proof}
Write \(p:\mathcal X'\to\mathcal X\) for the model morphism. Centers are
compatible with \(p\): for any \(x\in X^{\An}\),
\(c_{\mathcal X}(x)=p(c_{\mathcal X'}(x))\), because \(p\) is proper, so
the image of the center is a center, and centers are unique
(\Cref{prop:centers-proper-separated}).

\emph{Second identity.} If \(y\in\Sk(\mathcal X)\subset\Sk(\mathcal X')\),
then \(r_{\mathcal X'}(y)=y\) by \Cref{prop:skeleton-retraction} applied to
\(\mathcal X'\), whose retraction is the identity on its own skeleton.

\emph{Third identity.} \(r_{\mathcal X}\) takes values in
\(\Sk(\mathcal X)\subset\Sk(\mathcal X')\), on which \(r_{\mathcal X'}\)
is the identity by the previous paragraph.

\emph{First identity.} Set \(x'=r_{\mathcal X'}(x)\), and let
\(\xi=c_{\mathcal X}(x)\), \(\xi'=c_{\mathcal X'}(x)\), so
\(\xi=p(\xi')\). Let \(E_j\), \(j\in J(\xi)\), be the components of
\(\mathcal X_s\) through \(\xi\), with local equations \(z_j\). Since
\(p\) is an isomorphism on generic fibres, \(z_j\circ p\) is a local
equation at \(\xi'\) of the total transform \(p^*E_j\), a vertical divisor
on \(\mathcal X'\) supported on \(\mathcal X'_s\); in particular \(z_j\circ
p\) is a product of local equations of components of \(\mathcal X'_s\)
through \(\xi'\), times a unit. By the equality case of
\Cref{lem:retraction-minorant}, applied on \(\mathcal X'\) to that
monomial,
\[
 v_{x'}(z_j)=v_{r_{\mathcal X'}(x)}(z_j)=v_x(z_j)
 \qquad(j\in J(\xi)).
\]
Since \(r_{\mathcal X}\) depends on \(x\) only through the numbers
\(v_x(z_j)\) and the set \(J(\xi)\), it remains to see that
\(J(c_{\mathcal X}(x'))=J(\xi)\).  The two centers need not agree --- when
\(\xi\) is a closed point of \(E_i\) lying on no other component,
\(c_{\mathcal X}(x')\) is the generic point of \(E_i\), a strict
generization of \(\xi\) --- but \(c_{\mathcal X}(x')\) is always a
generization of \(\xi\) by \Cref{prop:skeleton-retraction}, so
\(J(c_{\mathcal X}(x'))\subseteq J(\xi)\).  For \(j\) in the smaller set
the displayed identity gives \(v_{x'}(z_j)=v_x(z_j)>0\), and
\(\sum_{j\in J(c_{\mathcal X}(x'))}b_jv_{x'}(z_j)=1
 =\sum_{j\in J(\xi)}b_jv_x(z_j)\) with all terms positive forces equality
of the two index sets.  We conclude
\(r_{\mathcal X}(x')=r_{\mathcal X}(x)\).
\end{proof}

The remaining assertion, the inclusion
\(\Sk(\mathcal X)\subset\Sk(\mathcal X')\) itself, is the substantial one.
It is a statement about how a monomial valuation for \(\mathcal X\) is
again monomial for the higher model, and it is proved by the toroidal
computation carried out for a stellar subdivision in
\Cref{ex:refinement-segment} and in general in
\cite[Proposition~3.4]{BFJ16}.

We can now settle the compatibility of integral affine structures announced
in \cref{sec:skeleton-affine}.

\begin{proposition}
\label{prop:affine-structure-refinement}
Let \(\mathcal X'\geq\mathcal X\) be SNC models, so that
\(\Sk(\mathcal X)\subset\Sk(\mathcal X')\). Then every cell of
\(\Delta_{\mathcal X}\) is a union of cells of \(\Delta_{\mathcal X'}\),
and for every cell \(\tau\) of \(\Delta_{\mathcal X'}\) contained in a cell
\(\sigma\) of \(\Delta_{\mathcal X}\), restriction maps
\(\operatorname{Aff}_{\mathbf Z}(\sigma)\) into
\(\operatorname{Aff}_{\mathbf Z}(\tau)\). Consequently the notions
\emph{rational affine}, \emph{rational piecewise affine}, and
\emph{rational polyhedral subdivision} are unambiguous on
\(\bigcup_{\mathcal X}\Sk(\mathcal X)\): they may be tested on any
sufficiently high model.
\end{proposition}

\begin{proof}
Let \(p:\mathcal X'\to\mathcal X\) and let \(h\in
\operatorname{Aff}_{\mathbf Z}(\sigma)\). By
\Cref{prop:integral-affine-intrinsic}\,(4) there is a vertical Cartier
divisor \(D\) on \(\mathcal X\), supported on \(\mathcal X_s\), with
\(h=\varphi_D|_\sigma\). Its pullback \(p^*D\) is a vertical Cartier divisor
on \(\mathcal X'\), and model functions are compatible with pullback:
\(\varphi_{p^*D}=\varphi_D\) as functions on \(X^{\An}\), by
\Cref{lem:model-function-well-defined}. Applying
\Cref{prop:integral-affine-intrinsic}\,(4) on \(\mathcal X'\), the
restriction of \(\varphi_{p^*D}\) to the cell \(\tau\) is integral affine
for \(\Delta_{\mathcal X'}\). Since \(\tau\subset\sigma\) and
\(\varphi_{p^*D}=\varphi_D\), this restriction is \(h|_\tau\).

That each cell of \(\Delta_{\mathcal X}\) is a union of cells of
\(\Delta_{\mathcal X'}\) follows from the inclusion of skeleta together
with the fact that both are finite polyhedral complexes and the cells of
\(\Delta_{\mathcal X'}\) meeting the relative interior of \(\sigma\) cover
it; concretely, it is exhibited by the toroidal charts of
\Cref{thm:skeleton-subdivision}. The final statement is then a formal
consequence: a subdivision or a piecewise affine structure witnessed on
\(\mathcal X\) is witnessed on \(\mathcal X'\), and any two models are
dominated by a third.
\end{proof}

\subsection{Subdivisions realized by blowups}

\begin{example}[Blowing up a stratum]
\label{ex:blowup-stratum}
Blowing up a stratum of an SNC special fiber performs a stellar subdivision
of the corresponding face. For the local model \(xy=\pi\), blowing up the
node inserts the valuation
\[
                         v(x)=v(y)=\tfrac12
\]
at the midpoint of the interval. Its exceptional component has multiplicity
two. Subsequent toroidal blowups insert rational points and produce rational
subdivisions.
\end{example}

\begin{example}[A subdivision computed in full: the segment]
\label{ex:refinement-segment}
We carry out \Cref{ex:blowup-stratum} completely, since it is the smallest
case in which every phenomenon of this section is visible.

\emph{The models.} Let \(\mathcal X\) be an SNC model containing the chart
\(\mathcal U=\Spec R[x,y]/(xy-\pi)\), with components \(E_x=\{x=0\}\) and
\(E_y=\{y=0\}\) of multiplicity one meeting in the node
\(P=\{x=y=0\}\). Its cell is the edge \(\sigma\) parametrized by
\(t=v(x)\in[0,1]\), with \(v(y)=1-t\), the vertices sitting at \(t=1\) and
\(t=0\). Let \(p:\mathcal X'\to\mathcal X\) be the blowup of \(P\).

\emph{The exceptional divisor.} On the \(x\)-chart write \(y=xy_1\); the
equation becomes \(x^2y_1=\pi\) with coordinates \(x,y_1\). The exceptional
divisor is \(F=\{x=0\}\), with \(\ord_F(x)=\ord_F(y)=1\) and
\(\ord_F(\pi)=2\), so \(b_F=2\) and the new vertex is \(\tfrac12\ord_F\),
sending \(x,y\mapsto\tfrac12\): the midpoint \(t=\tfrac12\), as
\Cref{ex:blowup-stratum} asserts. The special fibre near \(P\) becomes the
chain \(\widetilde E_x+2F+\widetilde E_y\), so \(\Delta_{\mathcal X'}\)
replaces \(\sigma\) by two edges meeting at the new vertex. As sets
\(\Sk(\mathcal X)=\Sk(\mathcal X')\) here: only the cell structure is finer.

\emph{The integral affine structures agree.} On \(\sigma_{\{x,F\}}\) the
coordinates satisfy \(\alpha_x+2\alpha_F=1\), and
\(t=v(x)=\alpha_x+\alpha_F=1-\alpha_F\), so \(\alpha_F=1-t\),
\(\alpha_x=2t-1\), and \(t\) runs over \([\tfrac12,1]\). Then
\[
 a_x\alpha_x+a_F\alpha_F=(2a_x-a_F)\,t+(a_F-a_x),
\]
and \((a_x,a_F)\mapsto(2a_x-a_F,\,a_F-a_x)\) is a bijection of
\(\mathbf Z^2\), with inverse \(a_x=c+d\), \(a_F=c+2d\). So the integral
affine functions of the subdivided edge are exactly \(ct+d\),
\(c,d\in\mathbf Z\), which is
\(\operatorname{Aff}_{\mathbf Z}(\sigma)|_{[1/2,1]}\) since on \(\sigma\)
one has \(a_x\alpha_x+a_y\alpha_y=(a_x-a_y)t+a_y\). This verifies
\Cref{prop:affine-structure-refinement} by hand, and shows that the
multiplicity \(b_F=2\) is exactly what keeps the structure unchanged: a new
vertex at \(\tfrac12\) with multiplicity one would double it. By
\Cref{rem:lattice-length} each new edge has lattice length
\(1/\operatorname{lcm}(1,2)=\tfrac12\), and the two sum to the lattice
length \(1\) of \(\sigma\).

\emph{Iterating.} Suppose adjacent components have multiplicities
\(b,b'\) and normalized vertices \(t,t'\).  Blowing up their node produces
an exceptional component of multiplicity \(b+b'\), whose normalized point is
\[
                 \frac{bt+b't'}{b+b'}.
\]
Thus blowing up the two nodes created in the first round inserts
\(t=\tfrac13\) and \(\tfrac23\), not dyadic quarter-points.  Iterated
toroidal node blowups insert Farey mediants, so that the vertex at \(p/q\) in
lowest terms carries multiplicity \(q\); by the Stern--Brocot construction
their union realizes every rational point of \(\sigma\).  Contrast
\Cref{ex:snc-chain-models}, where the model is \(xy=\pi^{\,n}\) and every
inserted vertex is integral of multiplicity one.  No blowup ever inserts an
irrational point: vertices of dual complexes are divisorial, hence rational.
Irrational points of \(\sigma\) lie on every skeleton containing \(\sigma\)
but are vertices of none.
\end{example}

\begin{example}[Barycentric subdivision]
\label{ex:refinement-triangle}
Return to \(\mathcal U=\Spec R[x,y,z]/(xyz-\pi)\) of
\Cref{ex:toric-triangle-chart}, with dual complex the standard triangle
\(\sigma\).

\emph{Blowing up the deepest stratum.} Let \(p:\mathcal U'\to\mathcal U\) be
the blowup of the closed point \(P=\{x=y=z=0\}\), with exceptional divisor
\(F\). Then \(\ord_F(x)=\ord_F(y)=\ord_F(z)=1\), so \(\ord_F(\pi)=3\) and
\(b_F=3\); the new vertex \(\tfrac13\ord_F\) sends each of \(x,y,z\) to
\(\tfrac13\), so it is the barycenter \(B\) of
\Cref{ex:monomial-triangle-points}. The new dual complex is the stellar
subdivision of \(\sigma\) at \(B\): three triangles \(BP_xP_y\),
\(BP_yP_z\), \(BP_zP_x\), so \(4\) vertices, \(6\) edges, \(3\) two-cells,
against \(3,3,1\) before.

\emph{Blowing up a curve stratum.} Blowing up \(\{x=y=0\}\) instead gives an
exceptional \(G\) with \(\ord_G(x)=\ord_G(y)=1\), \(\ord_G(z)=0\),
\(b_G=2\); the new vertex \(\tfrac12\ord_G\) is the midpoint
\((\tfrac12,\tfrac12,0)\) of the edge \(P_xP_y\), and the subdivision is the
stellar one along that edge. Blowing up the deepest stratum first and then
the strict transforms of the three one-dimensional strata realizes the
classical barycentric subdivision of \(\sigma\)
into six triangles, each of lattice area \(\tfrac1{12}\) in the coordinates
of \Cref{ex:toric-triangle-chart}, summing to \(\tfrac12\); by
\Cref{prop:affine-structure-refinement} the integral affine structure is
unchanged.

\emph{Nesting can be strict.} Unlike the one-dimensional case, blowing up a
closed point of \(\mathcal U\) that is not a torus fixed point --- say a
non-monomial point of the curve \(\{x=y=0\}\) --- creates a component whose
divisorial valuation is not monomial for \(\mathcal U\). Its vertex lies
outside \(\sigma\), and the new skeleton properly contains the old one. Both
phenomena, refinement of cells and genuine enlargement, occur in
\Cref{thm:skeleton-inverse-limit}, and only their combination makes the
union of all skeleta dense.
\end{example}

Not every abstract subdivision comes from one blowup, but projective
subdivisions do.

\begin{theorem}\label{thm:skeleton-subdivision}
Let \(\Delta'\) be a rational simplicial projective subdivision of
\(\Delta_{\mathcal X}\). There is a vertical blowup
\(\mathcal X'\to\mathcal X\) such that \(\mathcal X'\) is normal and
vertically \(\mathbf Q\)-factorial, its strata are indexed by the faces of
\(\Delta'\), and vertical Cartier divisors on \(\mathcal X'\) restrict to
affine functions on the faces of \(\Delta'\). After a further resolution,
one may work with an SNC model.
\end{theorem}
\begin{proof}
Because the subdivision is projective, it is defined by a strictly concave
rational piecewise affine support function $h$ on $\Delta_{\mathcal X}$ --- the
negative of a strictly convex one, the sign being fixed by the convention
that $v_\bullet(\mathcal I)$ is a \emph{minimum} of affine functions, so that
$\log|\mathcal I|$ is a maximum
(\Cref{lem:vertical-ideal-convex}\,(2)).  The maximal domains on which $h$ is
affine are exactly the maximal cells of $\Delta'$.  Multiplying $h$ by a positive integer makes every slope and
every value on a vertex integral.

Work first in an SNC chart
$\pi=u z_0^{b_0}\cdots z_r^{b_r}$.  An integral affine function on its
simplex is the valuation function of a vertical monomial
$\pi^a z_0^{a_0}\cdots z_r^{a_r}$.  Take the finite set of affine functions
whose \emph{minimum} is $h$ on this chart, and let $\mathcal I_h$ be the
fractional ideal generated by the corresponding monomials, so that
$v_\bullet(\mathcal I_h)=h$.  Adding a common affine
function multiplies all generators by one meromorphic monomial and does not
change the blowup.  Hence these local ideals agree on overlaps up to an
invertible factor and glue to a coherent vertical fractional ideal
$\mathcal I$ on $\mathcal X$.

Let $\mathcal X'$ be the normalization of
$\operatorname{Bl}_{\mathcal I}\mathcal X$.  On the chart where a generator
$m_0$ generates the transform of $\mathcal I$, its coordinate monoid is
obtained by adjoining all ratios $m/m_0$.  The inequalities
$v(m/m_0)\ge0$ say that $\ell_{m_0}\le\ell_m$ for every generator $m$, and so
cut out exactly the domain on which the affine function $\ell_{m_0}$ attached
to $m_0$ realizes the minimum, that is, the support function $h$.  Thus the normalized
blowup charts are the normal toroidal charts associated to the cells of
$\Delta'$.  Their toroidal strata correspond contravariantly to faces, and
the gluing is the face gluing of $\Delta'$.  This proves the asserted
incidence description.

A simplicial normal toroidal chart is $\mathbf Q$-factorial: each invariant
prime divisor becomes Cartier after multiplication by the index of the
primitive ray generators in the local lattice.  Since the subdivision is
simplicial, $\mathcal X'$ is therefore vertically $\mathbf Q$-factorial.
A vertical Cartier divisor is locally a monomial times a unit, so its
valuation function is affine on each cell.  Finally, in the standing
residue-characteristic-zero setting, resolve the toroidal singularities and
the total special fibre.  The resulting SNC model still dominates
$\mathcal X'$ and realizes a further subdivision.  This is the normalized
blowup construction of \cite[Theorem~3.15]{BFJ16}.
\end{proof}

\subsection{Cofinality of SNC models}

The residue-characteristic-zero hypothesis now enters decisively:
resolution of singularities makes SNC models cofinal among all models of the
smooth variety \(X\). It is therefore legitimate to use the directed system
of SNC models. Directed systems, rather than a chosen sequence, are
essential here.

\paragraph{The external input.}
We use exactly one theorem from outside this book at this point, and we
state its scope precisely.

\begin{quote}
\textbf{(R)} \emph{Resolution and principalization in characteristic zero.}
Let \(\mathcal Y\) be an integral scheme, proper and flat over \(R\), with
smooth generic fibre \(X\), and let \(\mathfrak b\) be a coherent ideal on
\(\mathcal Y\) cosupported on the special fibre. Then there is a projective
birational morphism \(\mathcal Y'\to\mathcal Y\), an isomorphism over the
regular locus of \(\mathcal Y\) where \(\mathcal Y_s\) is already SNC and
\(\mathfrak b\) already defines an SNC divisor, such that \(\mathcal Y'\) is
regular, \(\mathfrak b\mathcal O_{\mathcal Y'}\) is invertible, and the
union of \(\mathcal Y'_s\) with the exceptional locus and with the divisor
of \(\mathfrak b\mathcal O_{\mathcal Y'}\) is a simple normal crossings
divisor.
\end{quote}

We take \textup{(R)} as a relative principalization input.  The classical
sources are \cite{Hironaka1964I,Hironaka1964II};
\Cref{thm:appB-resolution} and \Cref{def:appB-log-resolution} record the
finite-type-over-a-field form, not a proof of the relative
finite-type-over-\(k[[t]]\) form quoted here. Two remarks
on hypotheses. First, our standing assumption that \(K\) is complete
discretely valued of residue characteristic zero makes \(R\) an
equicharacteristic complete discrete valuation ring, so \(R\simeq k[[t]]\)
by Cohen's structure theorem, with \(k\) of characteristic zero; the schemes
to which (R) is applied are of finite type over this excellent
\(\mathbf Q\)-algebra, which is the setting in which the statement is
standard, but the precise relative source remains part of this explicit
external input. Second, (R) is used \emph{only} in \Cref{prop:snc-cofinal} and
in the final step of \Cref{thm:skeleton-subdivision}; every other proof in
this chapter is characteristic free.

\begin{proposition}
\label{prop:snc-cofinal}
Assume the standing hypotheses and the input \textup{(R)}. Then:
\begin{enumerate}
\item every normal projective \(R\)-model of \(X\) is dominated by an SNC
  model;
\item any two SNC models are dominated by a third;
\item the SNC models form a directed subset of the set of models, cofinal in
  it; consequently
  \(\varprojlim_{\mathcal X\ \mathrm{SNC}}\Sk(\mathcal X)\) computes the
  same limit as any cofinal subsystem, and a model function is determined on
  some SNC model.
\end{enumerate}
\end{proposition}

\begin{proof}
(1) Let \(\mathcal Y\) be a normal projective model of \(X\). Apply
\textup{(R)} to \(\mathcal Y\) with \(\mathfrak b=\mathcal O_{\mathcal Y}\):
we obtain a projective birational \(\mathcal X\to\mathcal Y\) with
\(\mathcal X\) regular and \(\mathcal X_s\cup\mathrm{Exc}\) an SNC divisor;
in particular \(\mathcal X_s\) is SNC. As \(\mathcal X\to\mathcal Y\) is an
isomorphism over the generic fibre \(X\), which is smooth, \(\mathcal X\) is
again a model of \(X\), and it is projective over \(R\) because
\(\mathcal Y\) is and the morphism is projective.  To impose the harmless
indexing convention in the definition of an SNC model, take the barycentric
subdivision of the finite complex whose cells are indexed by irreducible
strata and realize it by the toroidal refinement of
\cite[Theorem~3.15]{BFJ16}.  On the resulting dominating SNC model, distinct
irreducible strata index distinct simplices. Hence \(\mathcal X\) is an SNC model dominating
\(\mathcal Y\).

(2) Let \(\mathcal X_1,\mathcal X_2\) be SNC models. By
\Cref{prop:common-domination} they are dominated by a common model
\(\mathcal Y\), which may be taken normal and projective; by (1),
\(\mathcal Y\) is dominated by an SNC model, which then dominates both
\(\mathcal X_i\).

(3) Directedness is (2), and cofinality is (1) together with the fact that
every model is dominated by a normal projective one --- take the
normalization of the graph closure, using
\Cref{prop:common-domination} again. For the last assertion, a model
function is by definition determined by a vertical
\(\mathbf Q\)-Cartier divisor on \emph{some} model \(\mathcal Y\); pulling
back to a dominating SNC model \(\mathcal X\) does not change the associated
function, by \Cref{lem:model-function-well-defined}, so the function is
determined on \(\mathcal X\).
\end{proof}

\begin{remark}[Where cofinality fails]
\label{rem:cofinality-scope}
\Cref{prop:snc-cofinal} is the exact point at which residue characteristic
zero is indispensable, and it should not be assumed silently elsewhere. In
mixed or positive residue characteristic no analogue is available in
arbitrary dimension; the substitute is de Jong's semistable alteration
theorem \Cref{thm:appB-de-jong}, which produces a generically finite cover
rather than a birational modification. The consequences for this chapter are
severe rather than cosmetic: without cofinality, the union of the skeleta of
those models that happen to be SNC need not be dense, and
\Cref{thm:skeleton-inverse-limit} may fail. This is the reason that the
Boucksom--Favre--Jonsson theory of \cref{part:bfj} carries a
residue-characteristic hypothesis in its foundational statements, and that
the trivially valued theory of \cref{part:bj}, which does not, has to be
built on a different foundation. The ledger of model-availability inputs in
\cref{chap:models} records the three regimes separately.
\end{remark}

\subsection{The inverse limit}

\begin{theorem}\label{thm:skeleton-inverse-limit}
Under the standing BFJ hypotheses, the map
\[
 X^{\mathrm{an}}\longrightarrow
 \varprojlim_{\mathcal X}\operatorname{Sk}(\mathcal X),
 \qquad
 x\longmapsto(r_{\mathcal X}(x))_{\mathcal X},
\]
where \(\mathcal X\) ranges over SNC models, is a homeomorphism.
Consequently,
\[
 \bigcup_{\mathcal X}\operatorname{Sk}(\mathcal X)
\]
is dense in \(X^{\mathrm{an}}\).
\end{theorem}

\begin{proof}
\emph{The map exists and is continuous.} By \Cref{prop:snc-cofinal} the SNC
models form a directed set, and by \Cref{prop:skeleton-nesting} the
retractions satisfy \(r_{\mathcal X}\circ r_{\mathcal X'}=r_{\mathcal X}\)
whenever \(\mathcal X'\ge\mathcal X\), so the skeleta with the transition
maps \(r_{\mathcal X}|_{\Sk(\mathcal X')}\) form an inverse system of
compact Hausdorff spaces and the family
\((r_{\mathcal X})_{\mathcal X}\) is compatible. Each \(r_{\mathcal X}\) is
continuous by \Cref{prop:skeleton-retraction}, so the induced map into the
inverse limit, which carries the subspace topology from the product, is
continuous.

\emph{Injectivity.} We use here the foundational separation input
\cite[Corollary~3.2]{BFJ16}: under the standing BFJ hypotheses, vertical
model functions on all models separate the points of \(X^{\An}\) and
generate its topology.  This input is essential; the divisor of a general
rational function has a horizontal part and does not supply an elementary
replacement.  Suppose \(x\neq y\), and choose a model function
\(\varphi\) with \(\varphi(x)\neq\varphi(y)\), and by
\Cref{prop:snc-cofinal}\,(3) an SNC model \(\mathcal X\) on which
\(\varphi=\varphi_D\) is determined. By
\Cref{lem:model-function-affine}, \(\varphi_D\circ r_{\mathcal X}
=\varphi_D\), so
\[
 \varphi_D(r_{\mathcal X}(x))=\varphi_D(x)
 \neq\varphi_D(y)=\varphi_D(r_{\mathcal X}(y)),
\]
whence \(r_{\mathcal X}(x)\neq r_{\mathcal X}(y)\).

\emph{Surjectivity.} Let \((y_{\mathcal X})_{\mathcal X}\) be a point of the
inverse limit. For each SNC model \(\mathcal X\), the set
\(F_{\mathcal X}=r_{\mathcal X}^{-1}(y_{\mathcal X})\) is a nonempty closed
subset of \(X^{\An}\): nonempty because \(y_{\mathcal X}\) itself lies in
it, and closed by continuity of \(r_{\mathcal X}\). The family
\((F_{\mathcal X})\) has the finite intersection property: given
\(\mathcal X_1,\ldots,\mathcal X_m\), choose by directedness an SNC model
\(\mathcal X'\) dominating them all; then
\(r_{\mathcal X_i}(y_{\mathcal X'})
 =r_{\mathcal X_i}(r_{\mathcal X'}(y_{\mathcal X'}))
 =y_{\mathcal X_i}\) using compatibility of the given family and
\Cref{prop:skeleton-nesting}, so \(y_{\mathcal X'}\in\bigcap_iF_{\mathcal
X_i}\). Since \(X^{\An}\) is compact (\Cref{thm:berk-proper-compact}), the
whole family has a common point \(x\), and by construction
\(r_{\mathcal X}(x)=y_{\mathcal X}\) for every \(\mathcal X\).

\emph{Homeomorphism and density.} The source is compact and the inverse
limit is Hausdorff, being a closed subspace of a product of compact
Hausdorff spaces; a continuous bijection between such spaces is a
homeomorphism. For density, let \(U\subset X^{\An}\) be nonempty open and
let \(x\in U\). Injectivity as proved above shows that the topology of
\(X^{\An}\) is generated by the model functions, so \(U\) contains a finite
intersection of sets \(\{|\varphi_{D_\ell}-\varphi_{D_\ell}(x)|<\epsilon\}\);
choosing an SNC model \(\mathcal X\) on which all the \(D_\ell\) are
determined and using \(\varphi_{D_\ell}\circ r_{\mathcal X}
=\varphi_{D_\ell}\), we get \(r_{\mathcal X}(x)\in U\cap\Sk(\mathcal X)\).
Hence every nonempty open set meets some skeleton.  The imported separation
and topology statement, as well as this inverse-limit conclusion in its
foundational form, is \cite[Corollary~3.2]{BFJ16}.
\end{proof}

\begin{definition}
A point lying on the skeleton of some SNC model is
\emph{quasimonomial}. A point corresponding to a rational point of a face
is divisorial on a suitable higher model.
\end{definition}

It follows that quasimonomial points are dense. In fact divisorial points are
already dense, but a whole face, not merely its rational points, is needed
for convexity and uniform estimates.

The two adjectives are sometimes used interchangeably in the literature, so
we separate them.

\begin{remark}[Quasimonomial versus divisorial]
\label{rem:quasimonomial-vs-divisorial}
\begin{enumerate}
\item A point \(\alpha\in\Delta_{\mathcal X}\) is divisorial if and only if
  its coordinates \(\alpha_j\), \(j\in J_+\), are rational. If they are, then
  \(r=1\) in \Cref{lem:skeleton-monomial-invariants} and the residue field
  has transcendence degree \(n\); conversely a divisorial valuation has
  value group \(c\mathbf Z\) with \(c\in\mathbf Q_{>0}\), since the group
  contains \(1=v(\pi)\), and its coordinates lie in that group. A rational
  point is realized as a vertex on a toroidal blowup by
  \Cref{thm:skeleton-subdivision}, applied to a subdivision having it as a
  vertex.
\item Divisorial points are therefore dense in \(X^{\An}\). Nevertheless the
  irrational ones cannot be discarded: they are needed for compactness, and
  for convexity arguments, which evaluate a function on a whole segment. A
  recurring error is to prove an estimate at all divisorial points and
  conclude it "by density", losing upper semicontinuity or a boundary value;
  the standing conventions on nets and regularized suprema exist to block
  this move.
\item Not every point of \(X^{\An}\) is quasimonomial. Type-IV points on a
  curve (\Cref{thm:berk-a1-classification}) lie on no skeleton, being limits
  of decreasing sequences of discs with empty intersection. The theorem says
  they are limits of quasimonomial points, not that they are quasimonomial.
\end{enumerate}
\end{remark}

\section{Vertical divisors and piecewise affine functions}
\label{sec:skeleton-model-functions}

Everything so far has produced spaces and maps. This section produces the
dictionary the rest of the book uses: an identification, in both directions,
between vertical divisorial data on models and piecewise affine functions on
skeleta. It is not a formality --- one direction reads a divisor off on a
cell, the other realizes prescribed polyhedral data by an algebraic blowup,
and the second is where \Cref{thm:skeleton-subdivision} is spent.

We now fix the sign convention that will remain in force throughout the
book.

\begin{definition}[Model function and sign convention]
Let \(D\) be a vertical Cartier divisor on a normal model \(\mathcal X\).
The model line bundle \(\mathcal O_{\mathcal X}(D)\) induces a metric on the
trivial bundle of \(X\). Its \emph{model function} is
\[
                 \varphi_D=\log|\mathcal O_{\mathcal X}(D)|.
\]
If \(D=\operatorname{div}(f)\) on a model chart, then \(f^{-1}\) is a frame
of \(\mathcal O_{\mathcal X}(D)\), and hence
\[
                    \varphi_D(x)=-\log|f(x)|=v_x(f).
\]
\end{definition}

There are two immediate sign tests.
\[
 D=\mathcal X_s=\operatorname{div}(\pi)
       \quad\Longrightarrow\quad \varphi_D\equiv v(\pi)=1.
\]
More generally, for every principal vertical divisor
\(D=\operatorname{div}(h)\),
\[
                         \varphi_D(x)=v_x(h).
\]
Thus effective vertical divisors give nonnegative model functions. These
tests rule out the opposite convention \(\log|f|\), which is also found in
the literature when the divisor or the relative metric is defined with the
opposite sign.

\subsection{Vertical divisors give affine functions}

\begin{lemma}
\label{lem:model-function-affine}
Let \(\mathcal X\) be an SNC model and let
\(D=\sum_{i\in I}a_iE_i\) be a vertical \(\mathbf Q\)-Cartier divisor
determined on \(\mathcal X\), that is, supported on \(\mathcal X_s\) with
\(a_i\in\mathbf Q\). Then on every cell \(\sigma_J\) of
\(\Delta_{\mathcal X}\),
\[
       \varphi_D(\iota_{\mathcal X}(\alpha))
                    =\sum_{j\in J}a_j\alpha_j.
\]
In particular \(\varphi_D\) is rational affine on every face, integral
affine when the \(a_i\) are integers, and
\[
                  \varphi_D\circ r_{\mathcal X}=\varphi_D
\]
as functions on \(X^{\An}\).
\end{lemma}

\begin{proof}
\emph{The formula on a cell.} Let \(\alpha\in\sigma_J\), with
\(J_+=\{j\in J:\alpha_j>0\}\), and let \(\xi_{J_+}\) be the center of
\(v_\alpha\) (\Cref{lem:skeleton-monomial-invariants}\,(1)). Near
\(\xi_{J_+}\) the only components of \(\mathcal X_s\) present are the
\(E_j\), \(j\in J_+\), so a local equation of \(D\) there is
\[
 f=u\prod_{j\in J_+}z_j^{a_j}
\]
with \(u\) a unit, in the sense that this identity holds after multiplying
by an integer clearing the denominators of the \(a_j\). Applying
\(v_\alpha\), which vanishes on units and sends \(z_j\) to \(\alpha_j\),
\[
 \varphi_D(v_\alpha)=v_\alpha(f)
 =\sum_{j\in J_+}a_j\alpha_j
 =\sum_{j\in J}a_j\alpha_j,
\]
the last equality because \(\alpha_j=0\) for \(j\in J\setminus J_+\). This
expression is manifestly rational affine in \(\alpha\), and integral affine
when all \(a_i\in\mathbf Z\), in the sense of
\Cref{def:integral-affine}.

\emph{The retraction identity.} Let \(x\in X^{\An}\) with center \(\xi\) on
\(\mathcal X\), and put \(\alpha_j=v_x(z_j)\) for \(j\in J(\xi)\). Near
\(\xi\) a local equation of \(D\) is \(f=u\prod_{j\in J(\xi)}z_j^{a_j}\)
with \(u\) a unit at \(\xi\), by the same argument. Hence
\[
 \varphi_D(x)=v_x(f)=\sum_{j\in J(\xi)}a_jv_x(z_j)
 =\sum_{j\in J(\xi)}a_j\alpha_j
 =\varphi_D(r_{\mathcal X}(x)),
\]
the last equality by the formula just proved applied to the point
\(\alpha=r_{\mathcal X}(x)\) of \(\sigma_{J(\xi)}\). This is the equality
case of \Cref{lem:retraction-minorant}: a vertical divisor is exactly a
monomial datum, and the retraction preserves monomial data.
\end{proof}

Reading the lemma as a statement about the whole space of vertical divisors
gives the clean linear-algebraic form of the correspondence in the affine
regime.

\begin{proposition}
\label{prop:vertical-divisor-affine-iso}
Let \(\mathcal X\) be an SNC model, let
\(\Div_0(\mathcal X)_{\mathbf Q}=\bigoplus_{i\in I}\mathbf Q\,E_i\) be the
space of vertical \(\mathbf Q\)-divisors supported on \(\mathcal X_s\), and
let
\[
 \mathcal A(\Delta_{\mathcal X})
 =\{h\in C^0(\Delta_{\mathcal X}):\ h|_{\sigma_J}
      \text{ is rational affine for every cell}\}.
\]
Then
\[
 \Div_0(\mathcal X)_{\mathbf Q}\longrightarrow
 \mathcal A(\Delta_{\mathcal X}),
 \qquad D\longmapsto\varphi_D|_{\Sk(\mathcal X)},
\]
is an isomorphism of \(\mathbf Q\)-vector spaces. Under it:
\begin{enumerate}
\item the vertex \(x_{E_i}\) has \(\varphi_D(x_{E_i})=a_i/b_i\), so
  \(D\) is recovered from \(h\) by \(a_i=b_i\,h(x_{E_i})\);
\item \(D\) is effective if and only if \(h\ge0\);
\item \(D=\mathcal X_s\) corresponds to \(h\equiv1\), and more generally
  \(D+c\,\mathcal X_s\) corresponds to \(h+c\);
\item \(D\) has integral coefficients if and only if \(h\) is integral
  affine on every cell, that is, if and only if \(b_ih(x_{E_i})\in\mathbf Z\)
  for every \(i\).
\end{enumerate}
\end{proposition}

\begin{proof}
The map is \(\mathbf Q\)-linear because \(D\mapsto\varphi_D\) is, both sides
being computed by the formula of \Cref{lem:model-function-affine}, and it
lands in \(\mathcal A(\Delta_{\mathcal X})\) by that lemma. Note that
\(\mathcal X\) is regular, so every \(\mathbf Q\)-divisor supported on
\(\mathcal X_s\) is \(\mathbf Q\)-Cartier and the map is defined on all of
\(\Div_0(\mathcal X)_{\mathbf Q}\).

\emph{Injectivity and (1).} Evaluate at the vertex \(x_{E_i}\), which is
the point of \(\sigma_{\{i\}}\) with \(\alpha_i=b_i^{-1}\); the formula
gives \(\varphi_D(x_{E_i})=a_ib_i^{-1}\). Hence \(h\) determines every
\(a_i\), and the map is injective.

\emph{Surjectivity.} Let \(h\in\mathcal A(\Delta_{\mathcal X})\) and set
\(a_i=b_ih(x_{E_i})\), \(D=\sum_ia_iE_i\). On a cell \(\sigma_J\), both
\(h\) and \(\varphi_D\) are affine and \(\sigma_J\) is a simplex with
vertices \(x_{E_j}\), \(j\in J\); an affine function on a simplex is
determined by its values at the vertices, and the two functions agree there
by construction. Hence \(h=\varphi_D\) on every cell, so everywhere.

\emph{(2)} If all \(a_i\ge0\) then \(\varphi_D\ge0\) by the formula. If some
\(a_{i_0}<0\) then \(h(x_{E_{i_0}})=a_{i_0}/b_{i_0}<0\).

\emph{(3)} For \(D=\mathcal X_s\) we have \(a_i=b_i\), and
\(\sum_jb_j\alpha_j=1\) is the defining equation of \(\sigma_J\); the second
statement follows by linearity.

\emph{(4)} By \Cref{def:integral-affine}, \(h|_{\sigma_J}\) is integral
affine exactly when it can be written \(\sum_ja'_j\alpha_j\) with
\(a'_j\in\mathbf Z\).  Evaluating at the vertex \(x_{E_j}\) gives
\(a'_j=b_jh(x_{E_j})=a_j\).  Thus such an integral representation exists
precisely when every \(a_j\) is integral, equivalently when
\(b_jh(x_{E_j})\in\mathbf Z\) for every \(j\).
\end{proof}

\begin{remark}[Variation of the model and variation by constants]
\label{rem:model-function-ambiguity}
\Cref{prop:vertical-divisor-affine-iso} says that on a \emph{fixed} SNC
model the correspondence is a bijection, with no ambiguity.  As models vary,
the same function is represented by pullbacks of a divisor; this is the
compatibility of \Cref{lem:model-function-well-defined}.  Multiples of the
special fibre are visible rather than ambiguous:
\(\varphi_{D+c\mathcal X_s}=\varphi_D+c\), so the corresponding metric is
scaled by \(\rme^{-c}\).  They become invisible only after one deliberately
quotients metrics by constants, or after one fixes a normalization such as
\(\varphi(x_0)=0\).  This is exactly
\Cref{prop:divisor-model-function-injective-mod-fibre}.
\end{remark}

\begin{example}[Model functions on the segment and on the triangle]
\label{ex:model-functions-computed}
\emph{On the segment.} Let \(\mathcal X\) contain the node chart of
\Cref{ex:refinement-segment}, with cell \(\sigma\) parametrized by
\(t=v(x)\in[0,1]\). For \(D=aE_x+bE_y\),
\[
 \varphi_D(v_t)=a\,t+b\,(1-t)=(a-b)t+b,
\]
an affine function with integral slope \(a-b\) and value \(b\) at \(t=0\).
The two sign tests read: \(D=\mathcal X_s=E_x+E_y\) gives
\(\varphi_D\equiv1\); \(D=\divis(x)=E_x\) --- valid in this chart, where
\(x\) has divisor \(E_x\) --- gives \(\varphi_D(v_t)=t=v_t(x)\), as the
convention demands.

After the blowup of \Cref{ex:refinement-segment}, with the extra component
\(F\) of multiplicity \(2\), a divisor \(D'=aE_x+cF+bE_y\) gives the
piecewise affine function with values \(a\) at \(t=1\), \(c/2\) at
\(t=\tfrac12\), and \(b\) at \(t=0\), interpolated affinely. Every
continuous function on \([0,1]\) that is affine on \([0,\tfrac12]\) and on
\([\tfrac12,1]\) arises this way, with \(c=2h(\tfrac12)\): the factor \(2\)
is the multiplicity, and forgetting it is the commonest computational error
in this circle of ideas.

\emph{On the triangle.} For the chart of \Cref{ex:toric-triangle-chart} and
\(D=a_xE_x+a_yE_y+a_zE_z\),
\[
 \varphi_D(v_\alpha)=a_x\alpha_x+a_y\alpha_y+a_z\alpha_z,
\]
which in the tropical coordinates \(u=(\alpha_x,\alpha_y)\) reads
\((a_x-a_z)u_1+(a_y-a_z)u_2+a_z\). Thus vertical divisors on this chart are
exactly the integral affine functions on the standard triangle, in
agreement with \Cref{ex:toric-triangle-chart}. After the barycentric
blowup of \Cref{ex:refinement-triangle}, with exceptional \(F\) of
multiplicity \(3\), the divisor \(D'=D+cF\) gives the function equal to
\(\varphi_D\) at the three original vertices and to \(c/3\) at the
barycenter, extended affinely on each of the three subtriangles. Taking
\(D=0\) and \(c=-1\), one obtains the function \(-\tfrac13\) at the
barycenter and \(0\) at the three vertices: the standard "tent function"
whose graph is a downward cone, which will reappear in
\cref{chap:bfj-model-forms} as the model function of a psh but not
pluriharmonic model metric.
\end{example}

\subsection{Piecewise affine functions come from vertical divisors}

A model function need not be determined on the chosen model. It becomes
affine after passing to a model on which its divisor is determined, and is
therefore piecewise affine on every lower skeleton.

\begin{proposition}
\label{prop:skeleton-model-pa}
Let \(\psi\) be a rational model function and let \(\mathcal X\) be an SNC
model. Then \(\psi|_{\operatorname{Sk}(\mathcal X)}\) is rational piecewise
affine. It is affine on every face if \(\psi\) is determined on
\(\mathcal X\).
\end{proposition}

\begin{proof}
By \Cref{prop:snc-cofinal}\,(3), choose an SNC model \(\mathcal X'\)
dominating \(\mathcal X\) on which
\(\psi\) is determined by a vertical \(\mathbf Q\)-Cartier divisor \(D'\).
By \Cref{lem:model-function-affine} applied on \(\mathcal X'\), the function
\(\varphi_{D'}\) is rational affine on each face of
\(\Delta_{\mathcal X'}\). By \Cref{prop:skeleton-nesting} we have
\(\Sk(\mathcal X)\subset\Sk(\mathcal X')\), and by
\Cref{prop:affine-structure-refinement} every cell of
\(\Delta_{\mathcal X}\) is a union of cells of \(\Delta_{\mathcal X'}\)
whose induced integral affine structure is the restricted one. Hence the
cells of \(\Delta_{\mathcal X'}\) inside a given cell \(\sigma\) of
\(\Delta_{\mathcal X}\) form a rational polyhedral subdivision of
\(\sigma\), and \(\psi|_\sigma=\varphi_{D'}|_\sigma\) is rational affine on
each piece. This is precisely the assertion that
\(\psi|_{\Sk(\mathcal X)}\) is rational piecewise affine. If \(\psi\) is
already determined on \(\mathcal X\), then
\Cref{lem:model-function-affine} applies directly on \(\mathcal X\) and no
subdivision is needed.
\end{proof}

\subsection{Vertical ideals and convexity}

Vertical ideals naturally produce convex rather than affine functions. If
\(\mathfrak a\) is a vertical fractional ideal and \(f_1,\ldots,f_N\)
generate it near the center of \(x\), put
\[
 \log|\mathfrak a|(x)
   =\max_\ell\log|f_\ell(x)|
   =-\min_\ell v_x(f_\ell).
\]
On a face this is the maximum of finitely many affine functions, hence is
convex and piecewise affine. The ultrametric inequality also gives
\[
               \log|\mathfrak a|(x)
        \leq \log|\mathfrak a|(r_{\mathcal X}(x)).
\]

Both assertions deserve proof, and both are now available.

\begin{lemma}
\label{lem:vertical-ideal-convex}
Let \(\mathcal X\) be an SNC model and \(\mathfrak a\subset
\mathcal O_{\mathcal X}\) a coherent ideal whose zero locus is contained in
\(\mathcal X_s\). Then:
\begin{enumerate}
\item \(\log|\mathfrak a|\) is well defined, independent of the chosen
  generators;
\item on every cell \(\sigma_J\) of \(\Delta_{\mathcal X}\), the function
  \(\alpha\mapsto\log|\mathfrak a|(v_\alpha)\) is a maximum of finitely many
  integral affine functions; in particular it is convex, continuous, and
  integral piecewise affine;
\item \(\log|\mathfrak a|\le\log|\mathfrak a|\circ r_{\mathcal X}\) on
  \(X^{\An}\);
\item if \(\mathfrak a=\mathcal O_{\mathcal X}(-D)\) is invertible with
  \(D\) vertical effective, then \(\log|\mathfrak a|=-\varphi_D\), which is
  affine on each cell. Convexity is therefore not a restriction but the
  precise measure of the failure of \(\mathfrak a\) to be invertible.
\end{enumerate}
The same statements hold for a fractional ideal, with "effective" dropped.
\end{lemma}

\begin{proof}
(1) If \((f_\ell)\) and \((g_m)\) both generate \(\mathfrak a\) near the
center of \(x\), each \(f_\ell\) is an \(\mathcal O_{\mathcal X,\xi}\)-linear
combination of the \(g_m\), so \(v_x(f_\ell)\ge\min_mv_x(g_m)\) by the
ultrametric inequality and nonnegativity of \(v_x\) on the local ring; by
symmetry the two minima agree.

(2) Fix a cell \(\sigma_J\) and let \(f_1,\ldots,f_N\) generate
\(\mathfrak a\) near \(\xi_J\). By \Cref{lem:monomial-newton}, each
\(\alpha\mapsto v_\alpha(f_\ell)\) is the minimum of the finitely many
integral affine functions \(\langle\cdot,\beta\rangle\) as \(\beta\) runs
over the vertices of the Newton polyhedron of \(f_\ell\). Hence
\[
 \log|\mathfrak a|(v_\alpha)
 =-\min_\ell v_\alpha(f_\ell)
 =\max_\ell\ \max_{\beta\in\mathrm{vert}\,N(f_\ell)}
   \bigl(-\langle\alpha,\beta\rangle\bigr),
\]
a maximum of finitely many integral affine functions. Note that the same
generators may be used at every point of \(\sigma_J\), because the center of
any \(v_\alpha\) with \(\alpha\in\sigma_J\) is a specialization-free
generalization of \(\xi_J\) inside \(E_J\); more precisely, it is
\(\xi_{J_+}\) with \(J_+\subset J\), and \(\xi_J\) specializes from
\(\xi_{J_+}\), so \(\mathcal O_{\mathcal X,\xi_J}\subset
\mathcal O_{\mathcal X,\xi_{J_+}}\) and the \(f_\ell\) still generate there.

(3) This is \Cref{cor:retraction-ideal-monotone}.

(4) If \(\mathfrak a=\mathcal O_{\mathcal X}(-D)\), then near any point a
single element \(f\) with \(\divis(f)=D\) locally generates, so
\(\log|\mathfrak a|=-v_\bullet(f)=-\varphi_D\), affine on each cell by
\Cref{lem:model-function-affine}.
\end{proof}

\begin{example}[Two vertical ideals on the segment]
\label{ex:vertical-ideal-segment}
\leavevmode\par\noindent
In the node chart \(\Spec R[x,y]/(xy-\pi)\), parametrize the cell by
\(t=v(x)\in[0,1]\).

\emph{A monomial ideal.} Let \(\mathfrak a=(x^p,y^q)\) with
\(p,q\in\mathbf N\). Then
\[
 \log|\mathfrak a|(v_t)=-\min\{pt,\ q(1-t)\}
 =\max\{-pt,\ qt-q\},
\]
a convex piecewise affine function with a single break at
\(t=q/(p+q)\), where the two affine pieces cross. Its two slopes are \(-p\)
and \(q\), both integers, and its total slope variation is \(p+q\). This
last number is the multiplicity that will reappear as a Monge--Amp\`ere mass
in \cref{chap:metrics}.

\emph{A non-monomial ideal.} Let \(\mathfrak a=(x-y,\ \pi)\). The Newton
polyhedron of \(x-y\) has vertices \((1,0)\) and \((0,1)\), so
\(v_t(x-y)=\min\{t,1-t\}\), and \(v_t(\pi)=1\); hence
\[
 \log|\mathfrak a|(v_t)=-\min\{\min\{t,1-t\},1\}
 =\max\{-t,\ t-1\},
\]
which is again convex piecewise affine with a break at \(t=\tfrac12\). Note
that a generator need not be a monomial for the conclusion to hold: this is
exactly the point of \Cref{lem:monomial-newton}, which converts an arbitrary
regular function into finitely many monomials for the purpose of computing
values on a skeleton.

\emph{The retraction inequality is strict.} Keep
\(\mathfrak a=(x-y,\pi)\) and let \(x_0\) be the closed point of the curve
\(\{xy=\pi\}\) with \(x=y=\pi^{1/2}\), a point with residue field
\(K(\pi^{1/2})\); it lies on the curve since \(xy=\pi\). At \(x_0\) one has
\(v_{x_0}(x)=v_{x_0}(y)=\tfrac12\), so \(r(x_0)\) is the midpoint
\(t=\tfrac12\), while \(v_{x_0}(x-y)=+\infty\) because \(x-y\) vanishes
there. Therefore
\[
 \log|\mathfrak a|(x_0)=-\min\{+\infty,\ 1\}=-1
 <-\tfrac12=\log|\mathfrak a|(r(x_0)),
\]
illustrating \Cref{lem:vertical-ideal-convex}\,(3). Retraction genuinely
loses information about non-monomial cancellation, and the inequality
records the direction in which it is lost.
\end{example}

\begin{theorem}\label{thm:skeleton-pl-realization}
Let \(h\) be a rational piecewise affine function on
\(\Delta_{\mathcal X}\), with slopes in the affine lattice described in
\Cref{sec:skeleton-snc}. Then there exist a model
\(\mathcal X'\) dominating \(\mathcal X\) and a vertical
\(\mathbf Q\)-Cartier divisor \(D\) on \(\mathcal X'\) such that
\[
                    \varphi_D|_{\operatorname{Sk}(\mathcal X)}=h.
\]
After multiplying \(h\) by a positive integer, \(D\) may be taken Cartier.
\end{theorem}

\begin{proof}
\emph{Step 1: reduction to an affine situation.} Let \(\Delta'\) be a
rational polyhedral subdivision of \(\Delta_{\mathcal X}\) on which \(h\) is
rational affine on each cell; such a subdivision exists by the definition of
a rational PL function, and after a further common refinement we may assume
it is simplicial. We may also assume it is projective: any rational
simplicial subdivision can be refined to a projective one, for instance by
taking the subdivision induced by the support function
\(\max_\ell\ell\) for a suitable finite family of rational affine functions
\(\ell\) --- concretely, by iterating stellar subdivisions at the vertices,
which is projective at each step and refines the given subdivision after
finitely many steps.

\emph{Step 2: realize the subdivision.} By
\Cref{thm:skeleton-subdivision}, there is a vertical blowup
\(\mathcal X'\to\mathcal X\) with \(\mathcal X'\) normal and vertically
\(\mathbf Q\)-factorial, whose strata are indexed by the faces of
\(\Delta'\) and on which vertical Cartier divisors restrict to affine
functions on the faces of \(\Delta'\). In particular the vertices of
\(\Delta'\) correspond to the components \(E'_k\) of \(\mathcal X'_s\), with
multiplicities \(b'_k\), and the identification is such that the vertex
attached to \(E'_k\) is the divisorial point \((b'_k)^{-1}\ord_{E'_k}\).

\emph{Step 3: read off the divisor.} Set
\[
 a'_k=b'_k\,h(\text{vertex attached to }E'_k)\in\mathbf Q,
 \qquad D=\sum_ka'_kE'_k .
\]
This is a vertical \(\mathbf Q\)-divisor, hence \(\mathbf Q\)-Cartier by the
vertical \(\mathbf Q\)-factoriality of \(\mathcal X'\). By the computation
of \Cref{lem:model-function-affine}, valid verbatim on the toroidal charts
of \(\mathcal X'\), the function \(\varphi_D\) is affine on each face of
\(\Delta'\) and takes the value \(a'_k/b'_k=h(\cdot)\) at each vertex. Since
an affine function on a simplex is determined by its vertex values, and
\(h\) is affine on each face of \(\Delta'\) by Step~1,
\[
 \varphi_D|_{\Sk(\mathcal X)}=h .
\]

\emph{Step 4: clearing denominators.} Let \(m\) be a positive integer with
\(ma'_k\in\mathbf Z\) for every \(k\). Then \(mD\) has integral coefficients
and, \(\mathcal X'\) being vertically \(\mathbf Q\)-factorial, a further
integral multiple of it is Cartier; the corresponding model function is
the corresponding multiple of \(h\).
\end{proof}

Combining the two directions gives the statement in the form used
throughout the rest of the book.

\begin{corollary}
\label{cor:pl-model-function}
Let \(\mathcal X\) be an SNC model. Then restriction to
\(\Sk(\mathcal X)\) induces a surjection
\[
 \{\text{rational model functions on }X^{\An}\}
 \longrightarrow
 \{\text{rational PL functions on }\Delta_{\mathcal X}\},
\]
whose kernel consists of the model functions vanishing on
\(\Sk(\mathcal X)\). Restricted to the model functions determined on
\(\mathcal X\), it is the isomorphism of
\Cref{prop:vertical-divisor-affine-iso} onto the rational affine functions.
Moreover the composite
\[
 \psi\longmapsto\psi|_{\Sk(\mathcal X)}\longmapsto
 \psi|_{\Sk(\mathcal X)}\circ r_{\mathcal X}
\]
is the identity on model functions determined on \(\mathcal X\), and is an
idempotent projection onto the rational model functions that factor through
\(r_{\mathcal X}\).
\end{corollary}

\begin{proof}
Surjectivity is \Cref{thm:skeleton-pl-realization} and well-definedness of
the map is \Cref{prop:skeleton-model-pa}. The description of the kernel is
tautological. The last assertion is
\Cref{lem:model-function-affine}: for \(\psi=\varphi_D\) determined on
\(\mathcal X\) one has \(\varphi_D\circ r_{\mathcal X}=\varphi_D\).  In
general \(\psi|_{\Sk(\mathcal X)}\) may have bends, so its realization can
require a higher model; nevertheless the composite factors through
\(r_{\mathcal X}\), has the same restriction to the skeleton, and applying
the composite again changes nothing.  Conversely every rational model
function factoring through \(r_{\mathcal X}\) is fixed.
\end{proof}

Thus vertical divisors and rational PL functions are two realizations of the
same object, provided models are allowed to vary.

\begin{remark}[What is not claimed]
\label{rem:pl-scope}
Three cautions, each of which will matter later.
\begin{enumerate}
\item The correspondence is between vertical divisors and PL functions, not
  between arbitrary divisors and arbitrary functions. A horizontal divisor
  does not give a model function; it gives a metric on a nontrivial line
  bundle, which is the subject of \cref{chap:metrics}.
\item Rationality is essential in \Cref{thm:skeleton-pl-realization}. A PL
  function on \(\Delta_{\mathcal X}\) with irrational breakpoints, or with
  irrational slopes, is not a model function. It is a uniform limit of
  model functions, and this is how such objects enter the theory; the
  completion of the space of model functions in the uniform norm is
  identified in \cref{chap:metrics}.
\item Convexity of a PL function on \(\Delta_{\mathcal X}\) is \emph{not}
  the correct notion of semipositivity. Convexity is a cellwise, coordinate
  dependent condition; semipositivity is defined by relative nefness of a
  model line bundle, and the two agree only in special situations such as
  the toric one of \cref{sec:toric-convex}. Conflating them is a
  well-known trap, and the book keeps the two words apart from the outset,
  as required by the conventions on positivity.
\end{enumerate}
\end{remark}

\section{Essential and canonical skeleta}
\label{sec:essential-skeleton}

The skeleton of a model depends on that model. Canonical subcomplexes arise
when the canonical bundle supplies a weight function. We record only the
definitions and scope needed later.

The word \emph{skeleton} is used for two kinds of object playing opposite
roles, and they should be separated at the outset.

A \emph{model skeleton} \(\Sk(\mathcal X)\) is a computational device: it
depends on the model, grows under refinement, is never canonical, and exists
so that the family of all of them exhausts \(X^{\An}\)
(\Cref{thm:skeleton-inverse-limit}). One chooses a model to make a
computation finite, as one chooses a triangulation to compute simplicial
cohomology; nothing attached to a single model skeleton is an invariant of
\(X\).

An \emph{essential} or \emph{canonical skeleton} is the opposite: a
distinguished compact subset of \(X^{\An}\) attached to \(X\) itself,
selected by a canonical differential form, birationally invariant, and
generally much smaller than any model skeleton. It does not grow under
refinement and exhausts nothing. Its interest is that limits of canonical
measures and of Ricci-flat metrics concentrate on it --- a theorem about
degenerations, not a tool for approximating analytic points.

The two coincide in the simplest examples, such as the Tate curve of
\Cref{ex:tate-degeneration}, but not in general: a rational surface has
arbitrarily large model skeleta and empty essential skeleton.

Let \(m>0\) and let
\[
              0\neq\omega\in H^0(X,mK_X).
\]
On an SNC model \(\mathcal X\), view \(\omega\) as a rational section of the
relative \(m\)-canonical bundle. At a divisorial point \(x_{E_i}\), define
\[
 \operatorname{wt}_{\omega}(x_{E_i})
   =\frac{\operatorname{ord}_{E_i}
      (\operatorname{div}_{\mathcal X}(\omega))+m}{b_i}.
\]
The added \(m\) is the logarithmic boundary contribution
\(m(\mathcal X_s)_{\mathrm{red}}\). Monomial interpolation extends the
weight function affinely on each face.

\begin{definition}
The \emph{Kontsevich--Soibelman skeleton} of \(\omega\) is the minimum locus
of \(\operatorname{wt}_\omega\). The \emph{essential skeleton} of \(X\) is
the union of these minimum loci over all \(m>0\) and all nonzero
\(m\)-pluricanonical forms. If there are no such forms, it is empty.
\end{definition}

Multiplying \(\omega\) by an element of \(K^\times\) adds a constant to its
weight function, so the minimum locus is unchanged. The model-independence
and birational invariance of these skeleta are deep theorems of
Kontsevich--Soibelman and Mustaţă--Nicaise. Under additional minimal-model
or Calabi--Yau hypotheses, the essential skeleton can be identified with the
dual complex of a dlt model and is sometimes called the canonical skeleton.
No such identification is assumed in the general pluripotential theory of
this book.

We record the elementary properties that do follow from the definition, so
that the reader can see exactly where the deep input begins.

\begin{lemma}
\label{lem:weight-elementary}
Fix \(m>0\) and \(0\neq\omega\in H^0(X,mK_X)\), and let \(\mathcal X\) be an
SNC model.
\begin{enumerate}
\item \(\operatorname{wt}_\omega\) is rational affine on every cell of
  \(\Delta_{\mathcal X}\); in particular it is continuous and attains its
  minimum on the compact set \(\Delta_{\mathcal X}\), so the
  Kontsevich--Soibelman skeleton is a nonempty closed subcomplex of some
  rational subdivision.
\item Replacing \(\omega\) by \(c\omega\) with \(c\in K^\times\) replaces
  \(\operatorname{wt}_\omega\) by \(\operatorname{wt}_\omega+v(c)\); the
  minimum locus is unchanged.
\item Replacing \(\omega\) by \(\omega^{\otimes k}\) replaces
  \(\operatorname{wt}_\omega\) by \(k\operatorname{wt}_\omega\); again the
  minimum locus is unchanged.
\item If \(\omega_1,\omega_2\) are nonzero sections of \(mK_X\) then
  \(\operatorname{wt}_{\omega_1+\omega_2}\ge
  \min\{\operatorname{wt}_{\omega_1},\operatorname{wt}_{\omega_2}\}\).
\end{enumerate}
\end{lemma}

\begin{proof}
(1) Write \(\divis_{\mathcal X}(\omega)=\sum_ic_iE_i+H\) with \(H\)
horizontal; along a cell only the vertical part contributes to the values on
divisorial points, and by construction
\(\operatorname{wt}_\omega(x_{E_i})=(c_i+m)/b_i\). By
\Cref{prop:vertical-divisor-affine-iso} the function on
\(\Delta_{\mathcal X}\) that is affine on every cell with these vertex
values is exactly \(\varphi_{D}\) for
\(D=\sum_i(c_i+m)E_i\), so it is rational affine on cells; the interpolation
prescribed after the displayed formula is precisely this affine extension.
Continuity and attainment of the minimum follow, and the minimum locus is
the union of the cells of a rational subdivision on which the affine
function achieves the minimum value.

(2) \(\divis_{\mathcal X}(c\omega)=\divis_{\mathcal X}(\omega)+v(c)
\mathcal X_s\) in the additive normalization, so
\(c_i\mapsto c_i+v(c)b_i\) and \(\operatorname{wt}\mapsto
\operatorname{wt}+v(c)\).

(3) The divisor of \(\omega^{\otimes k}\) is \(k\divis(\omega)\), viewed on
\(kmK_X\), and the boundary contribution \(m\) becomes \(km\); both terms
scale by \(k\).

(4) The order of vanishing of a sum is at least the minimum of the orders,
and the boundary contribution is the same for both.
\end{proof}

\begin{remark}
An arbitrary model skeleton is a computational stage for approximating all
analytic points. The essential skeleton is a birational invariant selected
by canonical forms. They answer different questions and should not be
identified.
\end{remark}

\begin{example}[Two extremes: the projective line and a Tate curve]
\label{ex:essential-extremes}
\emph{Empty essential skeleton.} Let \(X=\mathbf P^1_K\). Then
\(H^0(X,mK_X)=H^0(\mathbf P^1,\cO(-2m))=0\) for every \(m>0\), so the
essential skeleton is empty by definition. Nevertheless \(X\) has SNC models
with arbitrarily large dual complexes: the models
\(\widetilde{\mathcal C}^{(n)}\) of \Cref{ex:snc-chain-models} have dual
complex a segment subdivided into \(n\) unit edges. The same phenomenon
occurs for every rationally connected variety, and more generally whenever
the Kodaira dimension is \(-\infty\). One should therefore not think of the
essential skeleton as "the skeleton of \(X\)".

\emph{Maximal essential skeleton.} Let \(X\) be an elliptic curve with
multiplicative reduction as in \Cref{ex:tate-degeneration}. There
\(K_X=\cO_X\), a nonzero global \(1\)-form \(\omega\) exists, and the
weight function turns out to be constant, so the Kontsevich--Soibelman
skeleton is the entire dual complex, a circle. The essential skeleton is
therefore a circle of the same lattice circumference, and it is
\emph{equal} to the model skeleton of the minimal SNC model. This is the
prototype of a maximally degenerate Calabi--Yau degeneration.
\end{example}

\paragraph{Scope: what is used later, and what is not.}
\begin{enumerate}
\item \emph{Used essentially:} model skeleta, the retraction, the
  inverse-limit theorem, and the PL/vertical-divisor dictionary. Every
  theorem of \cref{part:bfj} and \cref{part:bj} mentioning a dual complex
  means a model skeleton.
\item \emph{Used only as a source of examples:} the essential skeleton, in
  \cref{chap:hybrid-spaces}, where limits of canonical measures concentrate
  on it, and in \cref{chap:mp-dual-complexes}, where the complex-analytic
  analogue is built. In both places the statements carry their own
  hypotheses.
\item \emph{Not used anywhere:} identification of the essential skeleton
  with the dual complex of a dlt minimal model, its birational invariance,
  and the contractibility results of the Kontsevich--Soibelman programme.
  These are true under stated hypotheses, but no proof here relies on them.
\end{enumerate}
A reader interested only in pluripotential theory may skip this section; a
reader who has met dual complexes through mirror symmetry should note that
the objects there are the essential skeleta, not the model skeleta that
carry the analysis.

\section{Three worked examples}
\label{sec:skeleton-worked}

The general theory is now complete. This section carries out three
computations from beginning to end: the local nodal model, source of every
one-dimensional statement in the book; a global degeneration of an elliptic
curve, the smallest example with nontrivial skeletal topology and the
running curve degeneration of \cref{sec:ch01-running-nodal} in final form;
and the two-dimensional toric chart of \Cref{ex:toric-triangle-chart}, on
which retractions, PL functions, and subdivisions are computed explicitly.

\subsection{The nodal degeneration revisited}

Consider
\[
                 \mathcal U=\operatorname{Spec}
                 R[x,y]/(xy-\pi)
\]
near its nodal special fiber. Its two components \(E_x=(x=0)\) and
\(E_y=(y=0)\) have multiplicity one. The skeleton is
\[
 \{(\alpha_x,\alpha_y)\in\mathbf R_{\geq0}^2:
                         \alpha_x+\alpha_y=1\}.
\]
Writing \(t=\alpha_x\), the associated valuation satisfies
\[
                         v_t(x)=t,\qquad v_t(y)=1-t.
\]
For \(D=aE_x+bE_y\),
\[
                         \varphi_D(v_t)=at+b(1-t).
\]
Blowing up the node inserts \(t=1/2\); suitable iterated blowups insert all
rational points. A vertical ideal generated by monomials
\(x^{p_\ell}y^{q_\ell}\) gives
\[
 \log|\mathfrak a|(v_t)
   =-\min_\ell\{p_\ell t+q_\ell(1-t)\},
\]
a convex piecewise affine function. This one-dimensional calculation is the
local model for every statement in
\Cref{sec:skeleton-model-functions}.

The retraction is equally explicit. A point \(w\) of the generic fibre, the
standard annulus of \Cref{ex:berk-standard-annulus}, has
\(r_{\mathcal U}(w)=v_t\) with \(t=v_w(x)\); the fibre over an interior
point \(t\in(0,1)\) is the set of points of the annulus with \(|x|=\rme^{-t}\),
which is an infinite union of open discs glued along the Gauss point
\(\zeta_{0,\rme^{-t}}\), and the fibre over an endpoint is the corresponding
residue class of the reduction map. This is \Cref{ex:retraction-annulus}.

\subsection{A semistable degeneration of an elliptic curve}
\label{sec:worked-tate}

\begin{example}[The Tate degeneration: circle, retraction, and weights]
\label{ex:tate-degeneration}
\index{Tate curve}\index{degeneration!elliptic curve}\index{skeleton!circle}
Assume, as throughout the chapter, that \(K\) is complete discretely valued
with \(v(\pi)=1\) and residue characteristic zero. Consider the plane
family
\[
 \mathcal E=\bigl\{XYZ=\pi\,(X^3+Y^3+Z^3)\bigr\}\subset\mathbf P^2_R .
\]

\emph{Step 1: the generic fibre is an elliptic curve.} Dividing by \(\pi\)
puts the equation of \(\mathcal E_K\) into Hesse form
\[
 X^3+Y^3+Z^3-3\lambda XYZ=0,\qquad \lambda=\tfrac1{3\pi},
\]
and a plane cubic in Hesse form is smooth exactly when
\(\lambda^3\neq1\) and \(\lambda\neq\infty\). Here
\(\lambda^3=(27\pi^3)^{-1}\), and \(\lambda^3=1\) would force
\(v(27\pi^3)=0\), that is \(3v(\pi)=v(27^{-1})=0\) since the residue
characteristic is zero and hence \(v(27)=0\); this contradicts
\(v(\pi)=1\). Therefore \(E:=\mathcal E_K\) is a smooth plane cubic, an
elliptic curve over \(K\).

\emph{Step 2: the special fibre is a triangle of lines.} Setting \(\pi=0\)
gives \(\mathcal E_s=\{XYZ=0\}\), the union of the three coordinate lines
\[
 E_X=\{X=0\},\qquad E_Y=\{Y=0\},\qquad E_Z=\{Z=0\},
\]
each a copy of \(\mathbf P^1_{\widetilde K}\), each of multiplicity one, so
the degeneration is strictly semistable if the total space is regular. The
three pairwise intersections are the coordinate points
\[
 E_Y\cap E_Z=[1:0:0],\quad E_Z\cap E_X=[0:1:0],\quad
 E_X\cap E_Y=[0:0:1],
\]
each a single reduced point, and \(E_X\cap E_Y\cap E_Z=\varnothing\).

\emph{Step 3: the total space is regular, with node charts.} Work near
\([1:0:0]\), in the affine chart \(X=1\) with coordinates \(y=Y/X\),
\(z=Z/X\). The equation reads
\[
 yz=\pi\,F(y,z),\qquad F(y,z)=1+y^3+z^3 .
\]
At the point \(y=z=0\) of the special fibre we have \(F(0,0)=1\), so \(F\)
is a unit on a neighbourhood; setting \(z'=z/F(y,z)\), which is a coordinate
change because \(F\) is a unit, the equation becomes
\[
                         yz'=\pi .
\]
This is the standard node chart of \Cref{ex:snc-chain-models}, so the local
ring at that point is regular and the special fibre is SNC there. Away from
the three coordinate points at most one of \(y,z\) vanishes, so the special
fibre is smooth there and the total space is regular by flatness. Hence
\(\mathcal E\) is an SNC model of \(E\), indeed strictly semistable.

\emph{Step 4: the dual complex is a circle.} The nonempty strata are the
three components and the three pairwise intersections; there is no triple
point. Hence \(\Sigma_{\mathcal E}\) consists of the three singletons and
the three pairs, and \(\Delta_{\mathcal E}\) is a triangle \emph{boundary}:
a cycle with three vertices and three edges, homeomorphic to a circle. This
is the case \(n=3\) of \Cref{ex:dual-complex-chain-cycle}. By
\Cref{rem:lattice-length}, all multiplicities being one, each edge has
lattice length \(1\), so the circle has total lattice circumference \(3\).

\emph{Step 5: the circumference is a valuative invariant.} The
\(j\)-invariant of the Hesse cubic is
\(j=27\lambda^3(\lambda^3+8)^3(\lambda^3-1)^{-3}\) with
\(\lambda^3=(27\pi^3)^{-1}\); substituting and simplifying,
\[
 j=\frac1{\pi^3}\cdot
   \Bigl(\frac{1+216\pi^3}{1-27\pi^3}\Bigr)^3 ,
\]
and the bracket is a unit of \(R\) since \(v(\pi)>0\). Hence
\(v(j)=-3\). For a curve with multiplicative reduction the Tate parameter
satisfies \(v(q)=-v(j)\), so \(v(q)=3\), matching Step~4. The general rule
is that the dual complex of a semistable degeneration of an elliptic curve
with multiplicative reduction is a circle of lattice circumference
\(v(q)=-v(j)\), equal to the number of components of the special fibre.
Replacing \(\pi\) by \(\pi^k\) multiplies the circumference by \(k\) after
resolving the resulting \(A_{k-1}\)-singularities, as in
\Cref{ex:snc-chain-models}.

\emph{Step 6: the retraction, in closed form.} Parametrize the circle by
lattice arclength \(\theta\in\mathbf R/3\mathbf Z\), the vertices sitting at
\(\theta=0,1,2\) and the edge from \(x_{E_X}\) to \(x_{E_Y}\) being
\([0,1]\). Let \(w\in E^{\An}\) with center \(\xi\). If \(\xi\) lies in the
smooth locus of a component, say \(E_X\), then \(J(\xi)=\{X\}\) and
\(r_{\mathcal E}(w)=x_{E_X}\): the whole preimage under reduction of that
smooth locus, an infinite union of residue classes, collapses to one point.
If \(\xi\) is the node \(E_X\cap E_Y\), then \(J(\xi)=\{X,Y\}\) and, in the
node coordinates \(u,u'\) with \(uu'=\pi\) of Step~3,
\[
   r_{\mathcal E}(w)=\theta=1-v_w(u)\in[0,1];
\]
the retraction reads the single number \(v_w(u)\) and forgets everything
else. The three node charts glue along the vertices into the circle, and the
resulting continuous surjection \(E^{\An}\to\mathbf R/3\mathbf Z\) is the
tropicalization of the Tate uniformization: under
\(E^{\An}\simeq\mathbf G_{m}^{\An}/q^{\mathbf Z}\), valid for split
multiplicative reduction, \(r_{\mathcal E}\) is induced by
\(w\mapsto v_w(s)\bmod v(q)\) for \(s\) the coordinate on
\(\mathbf G_m\). We use the plane model because it is elementary and
self-contained; the uniformization belongs to \cref{chap:cld-examples}.

\emph{Step 7: model functions.} By
\Cref{prop:vertical-divisor-affine-iso}, vertical \(\mathbf Q\)-divisors
\(D=a_XE_X+a_YE_Y+a_ZE_Z\) correspond bijectively to functions on the circle
affine on each edge, with \(\varphi_D(x_{E_X})=a_X\) and so on, all
multiplicities being one. Constants correspond to multiples of
\(\mathcal E_s\), consistent with \(\varphi_{\mathcal E_s}\equiv1\). And
there is no nonconstant \emph{convex} function on a circle: for a function
affine on edges, the outgoing slopes at all vertices telescope to zero, so
if all are nonnegative all vanish. This is the combinatorial shadow of the
maximum principle, reappearing in \cref{chap:model-calculations} as the
statement that a harmonic function on a compact metric graph is constant.

\emph{Step 8: the essential skeleton is the whole circle.} Since \(E\) is an
elliptic curve, \(K_E=\cO_E\) and there is a nowhere-vanishing global
\(1\)-form \(\omega\); concretely, for
\(G=XYZ-\pi(X^3+Y^3+Z^3)\), take the Poincaré residue
\(\omega=\operatorname{Res}(\Omega/G)\) with
\(\Omega=X\,\rd Y\wedge \rd Z+Y\,\rd Z\wedge \rd X+Z\,\rd X\wedge \rd Y\),
and take \(m=1\). The cyclic substitution \((X,Y,Z)\mapsto(Y,Z,X)\) is an
automorphism \(\varsigma\) of \(\mathbf P^2_R\) preserving \(G\), hence of
\(\mathcal E\) over \(R\); it fixes \(\Omega\), a cyclic permutation of
three coordinates having determinant \(+1\), so
\(\varsigma^*\omega=\omega\). As \(\varsigma\) permutes the three components
cyclically, the three vertex weights are equal, and by
\Cref{lem:weight-elementary}\,(1) the weight function, being affine on each
edge, is constant on the circle. Its minimum locus is therefore all of
\(\Delta_{\mathcal E}\): the Kontsevich--Soibelman skeleton of \(\omega\),
and \emph{a fortiori} the essential skeleton of \(E\), is the entire circle.
This is the coincidence promised in \Cref{ex:essential-extremes}; it is
special to maximally degenerate Calabi--Yau situations, and should be
contrasted with \(\mathbf P^1\), where the essential skeleton is empty while
model skeleta are arbitrarily large.

\emph{Step 9: refinement.} Blowing up one of the three nodes inserts a
component of multiplicity \(2\) at the midpoint of the corresponding edge
(\Cref{ex:refinement-segment}); the dual complex becomes a square, still a
circle of lattice circumference \(3\). No blowup changes the homotopy type
or the circumference: refinement subdivides but does not alter the intrinsic
metric structure of a one-dimensional skeleton, which is why the metric
graph of \cref{sec:metric-graphs} is well defined.
\end{example}

\subsection{A two-dimensional toric chart}
\label{sec:worked-toric}

\begin{example}[Retractions and PL functions on the triangle]
\label{ex:worked-triangle}
We complete the analysis of \(\mathcal U=\Spec R[x,y,z]/(xyz-\pi)\) begun in
\Cref{ex:toric-triangle-chart}. Recall \(\mathcal U_K\simeq
\mathbf G_{m,K}^2\) with coordinates \(x,y\), the dual complex is the
standard triangle \(\sigma\) with barycentric coordinates
\(\alpha=(\alpha_x,\alpha_y,\alpha_z)\), and tropical coordinates
\(u=(\alpha_x,\alpha_y)\) identify \(\sigma\) with
\(\Delta_{\mathrm{std}}\subset\mathbf R^2\).

\emph{The retraction is tropicalization.} For \(w\in\mathcal U_K^{\An}\) with
center on \(\mathcal U_s\),
\[
 r_{\mathcal U}(w)=\bigl(v_w(x),\,v_w(y),\,v_w(z)\bigr)
 =\bigl(v_w(x),\,v_w(y),\,1-v_w(x)-v_w(y)\bigr),
\]
so in tropical coordinates \(r_{\mathcal U}=\Trop\) restricted to the
preimage of \(\Delta_{\mathrm{std}}\), followed by the Gauss section
\eqref{eq:toric-gauss-section}. In other words, on a toric chart the
skeletal retraction \emph{is} the tropicalization map, and the skeleton is
the image of the Gauss section. This identification is the bridge between
this chapter and the toric calculus of \cref{sec:toric-convex}, and it is
the reason toric examples can be used to test every later normalization.

\emph{The fibres.} The fibre of \(\Trop\) over an interior point
\(u\in\Delta_{\mathrm{std}}\) is a two-dimensional analytic torus of
polyradius \((\rme^{-u_1},\rme^{-u_2})\); over a point of the open edge
\(\{\alpha_z=0\}\), the center of the retracted point moves to the generic
point of the curve \(E_x\cap E_y\) and the fibre is correspondingly larger.
Over a vertex, say \(P_x\), the fibre is the preimage under reduction of
the smooth locus of \(E_x\), a set with no finite-dimensional description at
all. The retraction is thus very far from being finite-to-one; what it
preserves is exactly the monomial data.

\emph{Model functions and their breaks.} By
\Cref{prop:vertical-divisor-affine-iso} and the coordinate computation of
\Cref{ex:toric-triangle-chart}, the vertical divisors on \(\mathcal U\)
correspond to the integral affine functions \(\langle m,u\rangle+c\) on
\(\Delta_{\mathrm{std}}\). To obtain a genuinely piecewise affine function
one must pass to a higher model. For instance
\[
 h(u)=\max\{0,\ -u_1,\ -u_2\}
\]
--- which on \(\Delta_{\mathrm{std}}\) equals \(0\), since
\(u_1,u_2\ge0\) there --- is uninteresting on this cell; a better test
function is
\[
 h(\alpha)=\min\{\alpha_x,\alpha_y,\alpha_z\},
\]
which is concave, piecewise affine, equal to \(0\) on the boundary of
\(\sigma\) and to \(\tfrac13\) at the barycenter, with the three domains of
affinity meeting along the medians. By
\Cref{thm:skeleton-pl-realization} it is the restriction of a model
function: indeed \(h=-\varphi_{D}\) with \(D=-F\) on the blowup
\(\mathcal U'\) of the closed point, where \(F\) is the exceptional divisor
of multiplicity \(3\); the vertex values are \(0\) at \(P_x,P_y,P_z\) and
\(\varphi_{-F}(B)=b_F^{-1}\ord_F(-F)=-\tfrac13\), so that
\(h(B)=-\varphi_{-F}(B)=\tfrac13\), and one checks that the three
subtriangles of \Cref{ex:refinement-triangle} are exactly the domains of
affinity of \(h\).

\emph{A vertical ideal.} Let \(\mathfrak a=(x,y,z)\subset
\mathcal O_{\mathcal U}\), the ideal of the closed point. Its zero locus is
contained in \(\mathcal U_s\), and by
\Cref{lem:vertical-ideal-convex},
\[
 \log|\mathfrak a|(v_\alpha)
 =-\min\{\alpha_x,\alpha_y,\alpha_z\}=-h(\alpha),
\]
convex, piecewise affine, integral, and equal to \(-\tfrac13\) at the
barycenter. Blowing up \(\mathfrak a\) is exactly the blowup of
\Cref{ex:refinement-triangle}, and the identity \(\log|\mathfrak a|
=-\varphi_F\) --- more precisely \(\log|\mathfrak a|=\varphi_{-F}\) on
\(\mathcal U'\) --- is the general principle that the normalized blowup of
an ideal turns \(\log|\mathfrak a|\) into a model function. This is the
mechanism used in \Cref{thm:skeleton-subdivision} and again in
\cref{chap:bfj-envelopes}.

\emph{An irrational point revisited.} The point \(\alpha\) of
\Cref{ex:monomial-triangle-points} has smallest coordinate
\(\alpha_x=(\sqrt2-1)/2\), so
\(\log|\mathfrak a|(v_\alpha)=-(\sqrt2-1)/2\). The value is irrational, as
it must be: model functions take values in \(\mathbf Q\) at divisorial
points but not at quasimonomial ones, and it is the continuity of PL
functions on the whole cell, rather than their values on rational points,
that makes the theory work.
\end{example}

\section{Dictionary established in this chapter}

\[
\begin{array}{ccl}
E_i\subset\mathcal X_s
 &\longleftrightarrow&
 b_i^{-1}\operatorname{ord}_{E_i}\text{, a vertex},\\[1mm]
E_J\neq\varnothing
 &\longleftrightarrow&
 \sigma_J\text{, a face of }\Delta_{\mathcal X},\\[1mm]
\alpha\in\sigma_J
 &\longleftrightarrow&
 v_\alpha\text{, a monomial valuation},\\[1mm]
D\text{ vertical Cartier}
 &\longleftrightarrow&
 \varphi_D\text{ affine on a determining skeleton},\\[1mm]
\mathfrak a\text{ vertical ideal}
 &\longleftrightarrow&
 \log|\mathfrak a|\text{ convex PL on each face}.
\end{array}
\]

In more detail, and with the precise references, the chapter establishes the
following correspondences. Each row is an equivalence unless the arrow is
one-directional.

\begin{center}
\begin{tabular}{@{}lll@{}}
\toprule
Model side & Polyhedral side & Reference\\
\midrule
component \(E_i\), multiplicity \(b_i\)
 & vertex, at \(b_i^{-1}\ord_{E_i}\)
 & \Cref{lem:skeleton-monomial-invariants}\\
stratum \(E_J\neq\varnothing\)
 & cell \(\sigma_J\), \(\dim=\#J-1\)
 & \Cref{lem:dual-complex-simplicial}\\
point \(\alpha\in\sigma_J\)
 & monomial valuation \(v_\alpha\)
 & \Cref{lem:monomial-valuation}\\
center of \(v_\alpha\)
 & smallest face containing \(\alpha\)
 & \Cref{lem:skeleton-monomial-invariants}\\
local equation of \(E_j\)
 & coordinate \(\alpha_j\)
 & \Cref{prop:integral-affine-intrinsic}\\
vertical divisor on \(\mathcal X\)
 & affine function on cells
 & \Cref{prop:vertical-divisor-affine-iso}\\
integral coefficients
 & integral affine structure
 & \Cref{def:integral-affine}\\
vertical ideal \(\mathfrak a\)
 & convex PL, \(\log|\mathfrak a|\)
 & \Cref{lem:vertical-ideal-convex}\\
\(f\) regular near \(\xi_J\)
 & concave PL, \(v_\bullet(f)\)
 & \Cref{lem:monomial-newton}\\
vertical blowup \(\mathcal X'\ge\mathcal X\)
 & subdivision, \(\Sk\) grows
 & \Cref{prop:skeleton-nesting}\\
rational projective subdivision
 & toroidal blowup realizing it
 & \Cref{thm:skeleton-subdivision}\\
rational PL function
 & vertical \(\mathbf Q\)-Cartier divisor
 & \Cref{thm:skeleton-pl-realization}\\
all SNC models
 & \(X^{\An}=\varprojlim\Sk(\mathcal X)\)
 & \Cref{thm:skeleton-inverse-limit}\\
pluricanonical form \(\omega\)
 & weight function, minimum locus
 & \Cref{lem:weight-elementary}\\
\bottomrule
\end{tabular}
\end{center}

Two of these deserve a warning label. The row "vertical blowup
\(\rightsquigarrow\) subdivision" is one-directional as stated: not every
subdivision comes from a blowup, only the projective ones, which is exactly
the content of \Cref{thm:skeleton-subdivision}. And the last row is not an
equivalence at all: the weight function determines a distinguished
subcomplex, but the essential skeleton does not determine \(\omega\), nor
does it determine any model.

\Cref{chap:metrics} uses the same vertical data to construct metrics on line bundles
and atomic measures from intersection numbers. \Cref{chap:bfj-model-forms}
upgrades the last two rows into a theory of closed \((1,1)\)-forms as
direct limits, and \cref{chap:bfj-psh} shows that psh functions restrict to
\emph{convex} functions on the cells of a skeleton, sharpening the row on
vertical ideals.

\section*{Source notes}

Skeleta of models go back to Berkovich's proof that smooth analytic spaces
are locally contractible, where the polystable case of
\Cref{thm:skeleton-deformation} is established; our reproduction follows his
construction, and the precise statement is his Theorem~8.1 in \emph{Smooth
\(p\)-adic analytic spaces are locally contractible}, Invent. Math.
\textbf{137} (1999). Thuillier developed the same picture for the special
fibre of a formal model and used it to prove birational invariance results
for dual complexes; his memoir is the source of the "\(\beth\)-space"
formalism used implicitly in our continuity proof. General references for
Berkovich geometry are \cite{Berk12} and \cite{Berk93}.

The elementary embedding, retraction, compatibility under domination, and
inverse-limit theorem are presented in the precise form used here in
\cite[Sections~3--4]{BFJ16}; the realization of subdivisions is the toroidal
theorem \cite[Theorem~3.15]{BFJ16}. The proofs above are self-contained
apart from two appeals to that source: the inclusion
\(\Sk(\mathcal X)\subset\Sk(\mathcal X')\) under domination, and the
toroidal construction realizing a projective subdivision. The
well-definedness of \(v_\alpha\), the Newton-polyhedron formula, the
intrinsic description of the integral affine structure, the
monomial-minorant inequality, and the divisor/PL dictionary on a fixed model
are proved here in full.

The only external input of substance is resolution of singularities in
characteristic zero, \cite{Hironaka1964I,Hironaka1964II}, used through
\Cref{thm:appB-resolution} for \Cref{prop:snc-cofinal}. We have isolated it
in the displayed hypothesis (R), so that a reader in positive or mixed
residue characteristic can see what is lost; the substitute there is de
Jong's alteration theorem \cite{deJong1996}, quoted as
\Cref{thm:appB-de-jong}, which is not birational and therefore does not
restore \Cref{thm:skeleton-inverse-limit}. See
\Cref{rem:cofinality-scope}.

The monomial-valuation calculus of \cref{sec:skeleton-monomial} is classical
valuation theory in modern dress, and the identification of quasimonomial
points with Abhyankar valuations in
\Cref{lem:skeleton-monomial-invariants} is standard; the form in which they
enter \cref{part:bfj} is that of \cite{BFJ15,BFJ16} and the survey
\cite{BFJ16b}.

Canonical weight functions and essential skeleta originate with
Kontsevich--Soibelman and were developed systematically by
Mustaţă--Nicaise and Nicaise--Xu. That material is included for orientation
and is not used as an unstated hypothesis later. Birational invariance of
the essential skeleton, its identification with the dual complex of a dlt
minimal model under Calabi--Yau hypotheses, and the resulting low-dimensional
topological classification lie outside our scope; a reader pursuing them
should note that those authors' normalizations of the weight function differ
from ours in the treatment of the logarithmic boundary term. Ours is fixed
by the two sign tests of \cref{sec:skeleton-model-functions}.

The toric computations follow the conventions of \cref{sec:toric-convex},
hence \cite{CLS11} for toric geometry and \cite{BGPS14} for the metric
normalizations; the identification of the skeletal retraction with
tropicalization in \Cref{ex:worked-triangle} is the elementary case of
\cref{chap:analytic-tropicalization}. The Hesse pencil computation in
\Cref{ex:tate-degeneration} is classical, and we prefer it to the Tate
uniformization because it needs nothing beyond the definition of an SNC
model.
